% Version 0.29. Actual support realization, with the original valued point.
% Source pin: 7cbf7a0c32fca1badf05931879a37c3f8a2b03b4.
\section{From a case-II valued point to its promodular support}
\label{sec:symplectic-promodular-realization}

\status{This chapter constructs the actual case-II witness required by
\texttt{ProModAlt}.  Its automorphic inputs are the literal
\texttt{SeedIIAlt} and \texttt{PMExAlt}, with the same fixed embedding
$\iota$.  The cohomology, residual vanishing and Ferrero--Washington
interfaces used below have already been reconstructed in BX, BR, AX
and TW.  In particular, an automorphic residual seed is not being
identified with the original valued point: the component and
coefficient arguments transporting support between them are part of
the proof.  The case-III boundary is recorded separately.}

\subsection{Objects that must remain fixed}

Use the original $\Omega/\mathbb Q$, odd prime $\ell$, coefficient
DVR $O\subset\overline{\mathbb Q}_\ell$, actual datum $D$, and all
of its decomposition groups and inertia groups.  Put $G_F=
\operatorname{Gal}(\Omega/F)$.  A case-II datum means the exact
\texttt{Ctx2}: $D$ is good and is the datum
\texttt{OfChar} for an original host character in Case II, with
the prescribed parameters $(n_{\rm II},1,1)$ of AX5.  In
particular its $\xi$ is a primitive odd dyadic character, and
the original auxiliary factor $\xi_0$ equals $1$.

Fix once and for all $\iota:\overline{\mathbb Q}_\ell\simeq
\mathbb C$.  An ambient field $(K,V,\{a_{O'}\})$ and a small
pair $(r,u)$ have exactly the meaning in IX.  We suppose that
$(r,u)$ is an actual \texttt{KPair} and that $r_K$ is
\texttt{Star}, namely absolutely irreducible or the specified
nonsplit block extension.  Every later change of lattice is
conjugation over this same $K$.

Write
\[
 \alpha=(1,0),\quad \beta=(-1,1),\quad
 \tau=(1,1),\quad \sigma=(0,1),\qquad
 \mathcal R=\{\pm\alpha,\pm\beta,\pm\tau,\pm\sigma\}.
 \tag{UX.1}
\]
The names $\tau,\sigma$ in this chapter denote roots, not Galois
elements.  They are source \texttt{wTh,wPs}.  The odd roots are
$\mathcal R_o=\{\pm\alpha,\pm\beta,\pm\tau\}$.  The matrix
weights, in the original basis $0,1,2,3$, are
\[
 (w_{ij})=
 \begin{pmatrix}
 0&\alpha&\sigma&\tau\\
 -\alpha&0&\beta&\sigma\\
 -\sigma&-\beta&0&\alpha\\
 -\tau&-\sigma&-\alpha&0
 \end{pmatrix}.
 \tag{UX.2}
\]
The three triangular patterns used below are
\[
 N_P=\{\beta,-\alpha,-\tau,-\sigma\},\quad
 N_M=\{-\alpha,-\beta,-\sigma,-\tau\},\quad
 Q=\{-\alpha,-\sigma,-\tau\},\quad
 V_0=\{\beta,-\alpha,-\tau\}.
 \tag{UX.3}
\]
Support in a root set means that every off-diagonal entry
whose weight is outside that set is zero.  An entry in a
root is \emph{present} if some such entry of some $r(g)$ is
nonzero.  This is the literal \texttt{Ent}; presence of an
entry of a local flag is stated separately.

A tower $T$ consists of the same label
$\mathrm{pr}\in\{0,1\}$, an imaginary quadratic field $K_0$,
a totally real abelian field $F$, their compositum $L=FK_0$,
a separating $g_0\in G_F$, and exactly one chosen place over
$\ell$ for each place of $F$, with the prescribed choice on
$K_0$.  Put $d=[F:\mathbb Q]$ and $s=\#\{v\mid\ell\}$.
The label permutes $(0,1,2,3)$ either trivially or by
$(2,3,0,1)$.  The relabelled characters are denoted $\chi_i$.
For a root $w=(a,b)$ put
\[
 \chi_w=(\chi_0/\chi_1)^a(\chi_0/\chi_2)^b.
 \tag{UX.4}
\]
The framed pair $P=(\rho,\{n_v\})$ has multiplier $D.\vartheta$,
$\rho(g_0)$ diagonal, and $n_v$ lower unipotent symplectic;
$n_v^{-1}\rho(G_{F_v})n_v$ is upper triangular.  Its local
characters are denoted $\psi_{v,i}$.  The frame kills the
specified inertia elements outside $\ell$ and a finite
auxiliary set $W$ whose original residual representation is
trivial.  It also imposes $\psi_{v,i}=1$ on the source's
$\mu_v$.  These are exactly $(T.\mathtt{fr}\,W).\mathtt{IsPair}$.
A residual pair has diagonal and ordinary characters
$\bar\chi_i$, is locally constant, and has either support
$N_P$ with flag support $Q$, or support $N_M$.

We retain all clauses of \texttt{Ok}, \texttt{Two},
\texttt{Fit}, \texttt{Deep} and \texttt{Le}.  In particular
$T\le T'$ means $K_0,\mathrm{pr}$ are unchanged,
$F\subset F'$, the quotient has $\ell$-power exponent, and
$g'_0$ restricts to an $\ell$-power of $g_0$.
\texttt{Deep} says $\psi|_{G_F}$ has order prime to $\ell$.
It does not merely assert that the layer index is large.

\subsection{The tower and the original lattice}

\begin{proposition}[UX1: the actual tame case-II tower and its extensions]
For every case-II $D$ and permitted label there exists an
actual tower $T_0$ satisfying \texttt{Ok}, \texttt{Two} and
\texttt{TameL}.  For every finite set $W$ of source
\texttt{Splits} primes it has an extension $T_1$ with the
same properties and with \texttt{Dies}$\,(\ell,w,F_1)$ for
every $w\in W$.  From any such tower, and any prescribed
$B\in\mathbb N$, there is a cyclotomic layer $T\ge T_1$
which is deep and has $[F:\mathbb Q]\ge B$.
\end{proposition}

\begin{proof}
Separate the $\ell$-primary part of the Dirichlet character
$\psi$ into its factor at $\ell$ and the factor away from
$\ell$.  Let $\kappa$ be the Galois character of the latter.
It has $\ell$-power order and is unramified at $\ell$.  On
$\ker c_\ell\cap\ker\kappa$, the character $\psi$ has order
prime to $\ell$.  On inertia away from $\ell$, if
$\bar\Psi=1$ and $\kappa=1$, then $\psi=1$: the remaining
factor has order prime to $\ell$ and reduction is injective
on roots of that order.  This proves every clause of the
source's \texttt{Kap}, without changing $\psi$ or its field.

Choose three new primes and their quadratic characters
$\varphi,\eta,\eta'$, respectively even, odd, odd.  Each
character is ramified at its own new prime only; $\varphi$
splits the finite set $D.S$, $\eta$ splits $D.S$ and the
prime of $\varphi$, and $\eta'$ splits all these and the
prime of $\eta$.  For example choose the prime in the
required class modulo $8\prod_{q\in D.S}q$, with the extra
conditions for previously chosen primes.  Quadratic
reciprocity and the Chinese remainder theorem specify a
nonempty reduced residue class, and the prime-in-progression
theorem supplies a prime.  The character of
$\mathbb Q(\sqrt{\pm p})$, with its sign and congruence
chosen so the fundamental discriminant has only that prime,
has exactly these properties.  Equivalently this is the
quadratic, degree-one specialization of Chebotarev.

Inside the original $G_{\mathbb Q}$ define
\[
 H=\{g:\bar\Psi(g)=\varphi(g),\quad
                 \eta(g)=\eta'(g),\quad\kappa(g)=1\}.
 \tag{UX.5}
\]
It is open, normal and contains the commutator.  Let $F$ be
its fixed field and $K_0$ the fixed field of $\ker\eta$.
Every complex conjugation belongs to $H$, because $\Psi$ is
even, $\kappa$ has odd order, and both $\eta,\eta'$ are odd.
Thus $F$ is totally real abelian and $K_0$ is imaginary
quadratic.  At a ramified place of either $\eta$ or $\eta'$,
the equality between them on $H$ kills both inertia values;
else both are already unramified.  Consequently $FK_0/F$
is unramified everywhere at finite places.  The prescribed
splitting of $D.S$ in $K_0$ gives the tower's \texttt{split}
clause and permits the required places over $\ell$.

The three auxiliary characters can be independently adjusted
by inertia at their three distinct ramified primes, while
leaving $\Psi,\xi,\omega,\kappa$ unchanged.  This proves
that the conditions in (UX.5) have their stated index; in
particular the index is even.  It also permits the following
three elements of $H\cap\ker\eta$:
\[
\begin{array}{c|ccc}
 &\bar\xi&\bar\Psi&\bar\omega\\ \hline
 a&-1&1&-1\\
 u&z&1&1\\
 t&z'&-1&1
\end{array}
\qquad z^4\ne1,\quad (z')^2\ne1.
 \tag{UX.6}
\]
Here $a$ starts from an odd power of complex conjugation;
$u$ starts from inertia at $2$; and $t$ starts from the
nontrivial prime-to-$\ell$ component of the Case-II character.
The latter component is nontrivial by AX1.  If the first
choice of $t$ has $(z')^2=1$, replace it by $tu$.
All powers used to kill $\kappa$ are odd and so preserve
the indicated dyadic nontrivialities.

We give the quantitative point in the construction of $u$
and its local analogue.  Write $\psi^{2^a b}=1$ with $b$
odd, as in AX1.  The prescribed $n_{\rm II}$ satisfies
$2^{n_{\rm II}}\ge\ell^2 2^{a+2}$.  If the fourth power
of $\xi(\operatorname{Frob}_\ell)^{2^a b\ell^j}$ were $1$,
the primitive dyadic estimate AX2 would give the strict
reverse inequality.  The chosen Frobenius has $c_\ell=1$,
hence $\omega=1$.  It therefore gives, at every place over
$\ell$, an element with
$\omega=\bar\Psi=\kappa=1$ and $\bar\xi^4\ne1$.
At $2$, a generator of the dyadic inertia quotient has the
full primitive character order; the same odd-power argument
gives $\bar\xi^4\ne1$ after the required factors are killed.
These are the source's \texttt{bigL} and \texttt{big2}, not
an unspecified large-image assertion.

On $G_F$ the residual $\sigma$-character is $\varphi$ (its
inverse is the same character), and the $\alpha$-character
is $\bar\xi$.  Thus it is quadratic, unramified everywhere,
trivial at $\ell$, and is cut out by the original real
quadratic field fixed by $\ker\varphi$.  That field splits
at $2$ and $\ell$.  These facts prove \texttt{Two}.
The rows of (UX.6) verify the remaining character clauses
of \texttt{Ok}: a root character cannot equal the
cyclotomic residual character (use $a$ for even first
coordinate and $u$ for odd first coordinate); every root
is nontrivial (use $u$ for odd roots and $t$ for
$\pm\sigma$); the $\alpha$-character has order greater
than four; and no nontrivial quadratic character trivial
on $\ker\omega$ permutes the four $\bar\chi_i$.
For the last assertion apply such a permutation to the
character $1$: evaluation at $u$ excludes $\xi$ and
$\xi\Psi$, and evaluation at $t$ excludes $\Psi$.
The four eigenvalues at $t$ are $(-z',-1,z',1)$, all
distinct, so choose $g_0$ above $t$.

For completeness the local \texttt{locL} alternative is
already supplied by the Frobenius just constructed.  Square
it to satisfy both quadratic equalities in (UX.5).  Its
$\omega$ and $\sigma$ values remain $1$, whereas its
$\alpha$ value is not $1$.  At $\ell$, $\xi$ and
$\varphi$ are unramified; inertia in $H$ therefore kills
the whole residual representation.  The two assertions
about $\psi$ follow from the first paragraph.  Finally
inertia at $\ell$ has prime-to-$\ell$ image in $G/H$,
since $\kappa$ is unramified there and all remaining
characters in (UX.5) have order prime to $\ell$.
This is literal \texttt{TameL}.

Now let $W$ be as stated.  Its primes are different from
$\ell$: complete splitting of $\bar r_0$ and $\omega$ at
$\ell$ would contradict AX6's separation at the prescribed
local Frobenius.  For each $w\in W$, take the compositum
of the cyclic degree-$\ell$ extensions of $\mathbb Q$
unramified outside $w$.  Global reciprocity identifies it
with the degree-at-most-$\ell$ subfield of
$\mathbb Q(\mu_w)$: the $w$-power principal units have
pro-$w$ order and contribute no map to $\mathbb F_\ell$.
This is a finite, totally real, abelian extension,
unramified at $\ell$.  Compositing with $F$ kills exactly
all the characters quantified in \texttt{Dies}.  Its
relative exponent divides $\ell$; take $g'_0$ above
$g_0^\ell$.  Distinct residual characters stay distinct
under this power.  Nontrivial prime-to-$\ell$ characters
and the other clauses of \texttt{Ok,Two} persist under an
$\ell$-extension: their restriction maps have
$\ell$-primary kernels and prime-to-$\ell$ images.
The fixed quadratic fields are unchanged.

For the final assertion set
\[
 G_{F_m}=G_{F_1}\cap
 \ker\bigl(c_\ell^{\ell-1}\bmod\ell^{m+1}\bigr).
 \tag{UX.7}
\]
These are exactly the source \texttt{IsLayer} fields.
The $\ell$-primary part of $\psi|_{G_{F_1}}$ vanishes
on $\ker c_\ell$, by \texttt{Ok.kc}, and hence factors
through the cyclotomic image.  It kills that image's
prime-to-$\ell$ Teichmuller subgroup.  Consequently
it is trivial wherever $c_\ell^{\ell-1}=1$.
By compactness it
vanishes on the intersection in (UX.7) for all sufficiently
large $m$.  More explicitly, the decreasing closed sets
on which that part is nontrivial and
$c_\ell^{\ell-1}\equiv1\pmod{\ell^{m+1}}$ have empty
intersection; one is already empty.  Thus these layers
are deep.  Tameness gives an element of $G_{F_1}$ whose
cyclotomic image has infinite order: raise the original
$s_0$ to the prime-to-$\ell$ power supplied by
\texttt{TameL}.  Its orders in the finite cyclotomic
quotients are unbounded, so $[F_m:\mathbb Q]\to\infty$.
Use $g_0^{\ell^m}$, enlarge $m$ if necessary, and choose
the original-field place representatives.  This proves
all the assertions with the original $K_0$ and label.
\end{proof}

\begin{proposition}[UX2: a good residual lattice for the original star point]
There is a permitted label, chosen from the original
ordinary flag as below.  Fix an initial tower $T_0$
supplied by UX1 for that label and its separating element
$g_0$.  There is a pair $P_0$ on the original
$G_{\mathbb Q}$ and original $V$, with a conjugacy
$r_K\simeq(P_0)_K$, such that $P_0$ is a small
\texttt{QPair} with the original inertia characters and its
residue has literal \texttt{LatShape}, normalized at this
$g_0$.  For every compatible fit extension of $T_0$
(whose separating element is the prescribed $\ell$-power
of $g_0$), the restricted residual pair, with the
transported flags, is good: it either has both
$-\sigma,\beta$ present, or is in $V_0$ with all
$n_v(2,0)\ne0$ and with $-\alpha$ present.
\end{proposition}

\begin{proof}
At the original $G_\ell$, AX6 separates the two residual
$\lambda$-blocks.  The first ordinary flag line is therefore
contained in the span of $e_0,e_2$ after reduction.  One of
$u_{00},u_{20}$ is a unit.  Use the trivial label in the first
case and the symplectic permutation $(2,3,0,1)$ in the
second.  In the second case the first-column eigenvector
equation, for every $g\in G_\ell$, gives
$\bar\chi_0(g)=\bar\chi_2(g)$; hence
$\bar\Psi|_{G_\ell}=1$.  This is exactly
\texttt{LabOk}, including its implication when the label
is nontrivial.  The same residual flag calculation makes
the first two leading minors units.  Symplecticity supplies
the remaining minor.

Now choose $T_0$ and $g_0$ by UX1 for this label.
Work initially in a complete DVR
defining the small pair.  Hensel's lemma splits the
characteristic polynomial of $r(g_0)$ into its four simple
residual roots.  The corresponding spectral idempotents
give an eigenbasis congruent to the original one; rescale
the paired eigenvectors to make its matrix symplectic.
This uses only units and, if necessary, a finite extension
of this defining DVR.  The ambient $K,V$ are not changed.
Conjugate the representation by this matrix and perform
the symplectic lower--upper decomposition of the ordinary
flag matrix.  Its lower factor $n$ is unique with diagonal
$1$.  Its residue has all off-diagonal entries zero except
possibly the root $-\sigma$ entries.  Multiplication by
the upper factor changes neither the ordinary characters
nor the inertia equalities.  The resulting representation
still has diagonal residual representation and diagonal
$r(g_0)$.  These constructions are finite algebraic
operations on small data, so smallness and continuity
over a complete finite-residue DVR persist.

There are nonzero generic entries in $\sigma$ and
$-\sigma$, and in at least one of
$\beta,-\alpha,-\tau$.  For an absolutely irreducible
representation this follows because its matrix span is
all of $M_4(K)$.  In the nonsplit block case, the full
$2\times2$ block span gives the first two assertions.
If all three remaining entries were zero, the symplectic
relations would make the other off-block entries zero as
well.  The conjugating matrix preserves the block
parabolic: it is assembled from spectral idempotents of
an element already in that parabolic and a block
permutation.  A block-preserving conjugacy splitting the
representation would give a unipotent off-block splitting
of the original representation (factor the conjugating
matrix into its diagonal-block and unipotent parts).
That contradicts the specified \texttt{BlockExt}.

Here is the lattice selection, including the finite
calculation behind it.  Let $\pi$ be a uniformizer of a
complete DVR defining the pair, and measure nonzero
entries by their integral $\pi$-orders.  For each root
$w$ include the orders of all representation entries of
that weight and, for the four negative roots, the orders
of the independent flag entries $n_{10},n_{21},n_{20},n_{30}$.
All are positive except possibly the order at $-\sigma$.
The set of integral pairs $(A,B)$ satisfying
\[
 w_1A+w_2B+2k\ge0
 \quad\hbox{for every recorded order $k$ at root $w$}
 \tag{UX.8}
\]
is nonempty.  Maximize $B$, then $A$.  The nonzero
$-\sigma$ entry bounds $B$ above; one of
$\beta,-\alpha,-\tau$ bounds $A$ above at fixed $B$.
Because these are integer bounds the maxima are attained.
Adjoin a square root $q$ of $\pi$ in the fixed $K$ and
conjugate by the symplectic diagonal torus with parameters
$(q^A,q^B)$.  Equation (UX.8) says exactly that all new
representation and flag entries belong to $V$.

Call a root tight if equality holds in (UX.8), so its
residual representation entry or independent flag entry
is nonzero.  The failed moves $(1,0),(0,1),(-1,1),(1,1)$
at a lexicographic maximum force respectively a tight
root in the following four sets:
\[
 \{-\alpha,\beta,-\tau\},\quad
 \{-\sigma,-\beta,-\tau\},\quad
 \{-\sigma,\alpha,-\beta\},\quad
 \{-\sigma,-\alpha,-\tau\}.
 \tag{UX.9}
\]
For the last two moves note first that the first two
conditions give $A+B$ even.  Thus a drop by $2$ in an
inequality of that parity must start at equality.
For every tight root other than $-\sigma$, positivity
of its original order gives $w_1A+w_2B<0$; tightness
of $-\sigma$ gives $B\ge0$.  Substitution of the eight
roots in (UX.1) into these sign inequalities and (UX.9)
gives exactly the following six possibilities:
\[
\begin{array}{c|c}
 \hbox{two roots that are tight}&\hbox{all tight roots lie in}\\ \hline
 \beta,-\sigma&N_P\\
 -\alpha,-\sigma&Q\\
 -\tau,-\sigma&\{-\tau,-\sigma\}\\
 -\alpha,-\beta&N_M\\
 -\tau,-\beta&\{-\tau,-\sigma,-\beta\}\\
 -\tau,\alpha&\{-\tau,-\sigma,-\beta,\alpha\}
\end{array}
 \tag{UX.10}
\]
For example $\beta$ tight means $B<A$, while
$-\beta$ tight means $A<B$; $\pm\tau$ have the signs
of $\pm(A+B)$ and $\pm\alpha$ those of $\pm A$.
These inequalities rule out any additional root in each
row and give (UX.10) without a compactness assertion about
an infinite-dimensional building.

The residual vanishing supplied by BR7 and AX4 now removes
the last five rows.  Precisely, for any odd root, a
locally constant normalized cocycle, unramified at the
specified inertia and locally a coboundary at $\ell$,
is zero, and its local coboundary parameter is zero.
At the other places in $D.T$, the local $H^1$ vanishes
by the separated residual Frobenius and the exact finite
local duality supplied by BX.  Thus AX4's original
\texttt{UnrEv}, also after any residue-field extension,
applies.  Its coboundary constant is zero because the
root character is not $1$ at $g_0$; the local parameter
is zero by AX6's local separation.

To see the two particular consequences needed here, write
the first-column ordinary identity as
\[
 r_{i0}(g)+r_{i1}(g)n_{10}+r_{i2}(g)n_{20}
       +r_{i3}(g)n_{30}=n_{i0}\bar\chi_0(g).
 \tag{UX.11}
\]
If $-\alpha$ is not tight, $r_{10}=r_{32}=n_{10}=0$;
the row $i=3$ makes $r_{30}/\bar\chi_0$ a locally
trivial odd-root cocycle.  Its vanishing gives absence
of $-\tau$.  If $\beta$ and $\sigma$ are not tight,
then $r_{12}=r_{13}=0$; the row $i=1$ similarly removes
$r_{10},n_{10}$ and symplecticity removes $r_{32}$.
Thus
\[
 -\alpha\notin\mathrm{Tight}\Rightarrow
 -\tau\notin\mathrm{Tight},\qquad
 \beta,\sigma\notin\mathrm{Tight}\Rightarrow
 -\alpha\notin\mathrm{Tight}.
 \tag{UX.12}
\]
Inspection of (UX.10) leaves its first row only.
Here $\beta$ can occur only in the representation, and
the flag has support $Q$.  If $-\sigma$ is absent from
the representation, tightness forces $n_{20}\ne0$.
Equation (UX.11) at $i=2$ then gives
$\bar\Psi|_{G_\ell}=1$, and the $\beta$ entry cannot
be locally a coboundary, by the same odd-root vanishing.
These are exactly all clauses of \texttt{LatShape}.

Finally restrict to $G_F$.  Nonzero scalar cocycles in
the indicated roots stay nonzero.  Indeed, if a
normalized scalar cocycle vanished on $G_F$, evaluate
its cocycle identity at $g g_0g^{-1}$: since the two
diagonal characters at $g_0$ are distinct, its value
at every $g$ would be zero.  This proves the assertion
directly, even when $[F:\mathbb Q]$ is divisible by
$\ell$.  The transported lower flags satisfy
\[
 n'_{20}\bar\chi_0(g)=r_{20}(g)+\bar\chi_2(g)n_{20},
 \quad
 n'_{10}\bar\chi_0(g)=r_{10}(g)+\bar\chi_1(g)n_{10}
                                      +r_{12}(g)n_{20}.
 \tag{UX.13}
\]
They also have $n'_{21}=0$.  These identities follow by
equating the first two columns in $r(g)n=n'b$ with $b$
upper triangular.  If $r_{20}$ is present, it and the
$\beta$ entry remain present.  If it is absent,
(UX.13) keeps every $n'_{20}$ nonzero.  Were the
$-\alpha$ entry also absent, its second equation would
make the original $\beta$ entry locally a coboundary,
contrary to the preceding paragraph.  This gives the
two asserted good alternatives and completes the proof.
\end{proof}

\subsection{Finite cohomology and the three residual lines}

\begin{proposition}[UX3: simultaneous nonzero cup extensions with removable ramification]
Let $T$ be fit, let $k$ be any field receiving the original
residue field, and let finitely many locally constant
scalar cocycles $y_j$ have character $\chi_{-\sigma}$,
vanish at $g_0$, and be locally coboundaries at the chosen
places over $\ell$.  Allow ramification at a finite set
$W$ of source \texttt{Splits} primes.  After adding finitely
many further such primes, there are a common normalized
$\chi_{-\alpha}$-cocycle $x$, local constants $m_v$, and
cochains $z_j$ such that
\[
\begin{aligned}
 x(gh)&=x(g)+\chi_{-\alpha}(g)x(h),\quad x(g_0)=0,\\
 x(s)&=m_v-\chi_{-\alpha}(s)m_v\quad(s\in G_{F_v}),\\
 z_j(gh)&=z_j(g)+\chi_{-\tau}(g)z_j(h)
                 +y_j(g)\chi_{-\alpha}(g)x(h),\\
 z_j(s)&=o_{jv}-\chi_{-\tau}(s)o_{jv}
                 -y_j(s)\chi_{-\alpha}(s)m_v,
 \qquad z_j(g_0)=0.
\end{aligned}
 \tag{UX.14}
\]
There is a $g$ with $\chi_{-\alpha}(g)=1$ and $x(g)\ne0$.
All these cochains kill the enlarged frame's forbidden
inertia.  They produce actual residual symplectic pairs
with support $Q$, a common $-\alpha$ entry and common
$n_{v,10}$, while retaining the given $y_j$ and
$n_{v,20}$.  A tower extension supplied by UX1 removes
all the added ramification, including when $\ell=3$.
\end{proposition}

\begin{proof}
First we prove the finite-module assertion used in this
construction.  Let $M$ be a finite $\mathbb F_\ell[G_F]$-
module, unramified outside a finite $S_0$ containing the
places over $\ell$, and choose $0\ne m_0\in M$.
Chebotarev supplies infinitely many rational primes
splitting completely in the finite Galois extension
cut out by $M$, $F$, $\bar r_0$ and $\mu_\ell$.
Use their degree-one places in $F$ and remove $S_0$.
These primes satisfy exactly \texttt{Splits}; at them
the module and the mod-$\ell$ cyclotomic character are
trivial on the whole decomposition group.

At an added place impose the following local condition:
the restriction of a cocycle to inertia lies in
$\mathbb F_\ell m_0$.  Unramified cocycles give a copy
of $M$; an independent tame cocycle with image in
$\mathbb F_\ell m_0$ gives another factor of $\ell$.
Indeed the tame local relation is
$\operatorname{Fr}\tau\operatorname{Fr}^{-1}=\tau^{q_v}$,
and $q_v\equiv1\pmod\ell$.  The two homomorphisms from
the pro-$\ell$ quotient, one on Frobenius and one on
tame inertia, are independent.  The local condition
therefore has size $\ell|M|\ge2|M|$.  At $S_0$ impose
the zero cohomology class.

The exact finite Greenberg--Wiles identity supplied by
DX3, or its direct Poitou--Tate cardinal proof, now gives
\[
 2^r\le |\operatorname{Sel}_{S_0\cup Q_r}(M)|\,
       |M^*(1)|\,|M|^{[F:\mathbb Q]+|S_0|}.
 \tag{UX.15}
\]
Here $Q_r$ is any set of $r$ of the new places.  To
check the coarsening, bound each local $H^0$ by $|M|$,
each real-place contribution by $|M|$, and the dual
global invariants by $|M^*(1)|$; the factors $|M|^r$
from the new local conditions cancel those new local
$H^0$ bounds.  No constant on the right depends on $r$.
The group of actual cocycles unramified outside $S_0$
is finite by BX8's finite global cohomology and the
finite group of coboundaries; call its cardinality $N$.
Put $C=|M^*(1)|\,|M|^{[F:\mathbb Q]+|S_0|}$ and take
$r=NC$.  If every Selmer cocycle had zero inertia at
the added places, choosing one representative per
class would inject this Selmer group into those $N$
cocycles.  Equation (UX.15) would imply $2^{NC}\le NC$,
a contradiction.  Thus some cocycle $c$ has a nonzero
inertia value in $\mathbb F_\ell m_0$; the acting
element has trivial action on $M$.  This cocycle is
locally a coboundary at every old place.

We apply this assertion to an explicit finite module,
so no finite-dimensionality assumption on $k$ is hidden.
The image of the coefficient residue map is a finite
field $k_0$.  Each locally constant $y_j$ has finite
image because $G_F$ is compact.  Let $V_0'$ be their
finite-dimensional $k_0$-span in $k$ and put
$M=k_0\oplus\bigoplus_j V_0'$.  With
$\rho=\chi_{-\alpha}$ and $\sigma'=\chi_{-\sigma}$,
define
\[
 g\cdot(a,(b_j))=
  \bigl(\rho(g)a,
    (\rho(g)\sigma'(g)b_j+\rho(g)a\,y_j(g))_j\bigr),
 \qquad m_0=(1,0).
 \tag{UX.16}
\]
The cocycle identities for $y_j$ show directly that
this is an action.  It factors through a finite quotient
and kills the original forbidden inertia.  Project the
cocycle $c$ just constructed to its coordinates.  Its
first coordinate is $x$, the remaining coordinates are
$z_j$, and the cocycle identity is exactly (UX.14).
Writing $c(s)=g m-m$ at an old local group gives the
local identities in (UX.14), with the negatives of
the coordinates of $m$.  On the nonzero inertia value
in $\mathbb F_\ell m_0$, the first coordinate is
nonzero; the first part of (UX.16) then forces
$\rho(g)=1$.

Normalize at $g_0$ without losing this nonzero value.
Put $b=x(g_0)/(1-\rho(g_0))$ and subtract the
coboundary $b-\rho b$ from $x$.  For each $j$ put
$c_j=z_j(g_0)/(1-\rho(g_0)\sigma'(g_0))$ and replace
\[
 z_j(g)\longmapsto z_j(g)-
 \{c_j-\rho(g)\sigma'(g)c_j-y_j(g)\rho(g)b\},
 \qquad m_v\longmapsto m_v-b.
 \tag{UX.17}
\]
The denominators are nonzero by the separating $g_0$.
Substitution into (UX.14) proves the normalized
identities.  Since $\rho(g)=1$ at the witnessed
nonzero value, that value of $x$ is unchanged.

Here are the resulting matrices.  Write
\[
 U(x,y,t)=
 \begin{pmatrix}
 1&0&0&0\\x&1&0&0\\y&0&1&0\\t&y&-x&1
 \end{pmatrix},\qquad
 U(x,y,t)U(x',y',t')
 =U(x+x',y+y',t+t'+yx'-xy').
 \tag{UX.18}
\]
For the $j$th datum define
\[
 r_j(g)=U(x(g),y_j(g),2z_j(g)-x(g)y_j(g))
                         \operatorname{diag}(\bar\chi_i(g)),
 \quad
 n_{jv}=U(m_v,n_{jv}^{\rm old},2o_{jv}-m_vn_{jv}^{\rm old}),
 \tag{UX.19}
\]
where $y_j(s)=n_{jv}^{\rm old}
(1-\sigma'(s))$.  The diagonal conjugation scales
$x,y,t$ by $\rho,\sigma',\rho\sigma'$.
Equations (UX.14) and (UX.18) show both that $r_j$ is
a representation and that $n_{jv}^{-1}r_j(s)n_{jv}$
is the prescribed diagonal matrix.  They also give
the symplectic multiplier, normalization at $g_0$,
support $Q$, and every required inertia equality.
The common $x,m_v$ give literal \texttt{SameAl}; the
unchanged $y_j,n_{jv}^{\rm old}$ give literal
\texttt{SamePs}.  The nonzero $x$ gives presence of
$-\alpha$.  This verifies the simultaneous
\texttt{gysinAl} conclusion entry by entry.

Finally apply UX1 to the added rational primes.  At
one of them $w$, the residual diagonal characters are
trivial.  A matrix supported in $Q$ is unipotent of
nilpotence class at most two and has exponent $\ell$
for every odd $\ell$: (UX.18) gives $U^n=U(nx,ny,nt)$.
The pro-$\ell$ quotient of tame inertia is cyclic;
its unique nontrivial degree-$\ell$ character is the
one killed by the degree-$\ell$ cyclotomic subfield
used in \texttt{Dies}.  Wild inertia is pro-$w$ and
acts trivially.  Hence the restricted representation
is unramified there.  This proves the claim also at
$\ell=3$, without using the false assertion that
every $4\times4$ unipotent matrix in characteristic
three has exponent three.  Existing residual pairs
already unramified at $w$ need no such assertion.
\end{proof}

\begin{proposition}[UX4: the complete three-line case-II residual bridge]
For every tame \texttt{Ok,Two} tower $T_0$ and every
original \texttt{LatShape} residual pair, there is a
tame extension $T_1\ge T_0$ such that, for every deep
fit extension $T\ge T_1$, literal \texttt{SeedIIAlt}
at the fixed $\iota$ supplies an integral seed $Q$
whose residue is connected to every compatible
restriction of that original residual pair by a
chain of three polynomial residual pairs.  Every
endpoint in the chain is good in the literal sense.
\end{proposition}

\begin{proof}
Restrict the original pair and transport its ordinary
flags as in UX2; call the result $p$.  It has support
$N_P$ and flag support $Q$.  Its normalized entry
$y(g)=p_{20}(g)/\bar\chi_0(g)$ is a
$\chi_{-\sigma}$-cocycle, with local constant
$n_v=p_{v,20}$.  Apply UX3 (one may use the same
datum twice) to produce a residual pair $q$ supported
in $Q$, with $\mathtt{SamePs}(p,q)$ and a nonzero
$-\alpha$ entry.  Let $T_1$ kill its finitely many
new ramified primes.  This choice precedes $T$, the
seed and its coefficient ring.

Fix any later deep fit $T$ and use SeedIIAlt there.
It produces an actual integral seed whose residue $z$
is in $V_0$, with all $z_{v,20}\ne0$, and satisfies
\[
 \operatorname{Ent}_{-\alpha}(z)\quad\hbox{or}\quad
 \bigl(\operatorname{Ent}_{\beta}(z)\ \hbox{and}\
       \operatorname{Ent}_{-\tau}(z)\bigr).
 \tag{UX.20}
\]
Map its finite residue field into the given ambient
residue field; the compatible coefficient maps specify
this embedding.  Restrict $p,q$ to $T$, choosing their
flags over the tower's actual place representatives.
The identities (UX.13) transport \texttt{SamePs},
and the nonzero normalized $-\alpha$ cocycle survives
restriction: its character at the appropriate
$\ell$-power of $g_0$ is still nontrivial.  Auxiliary
inertia is killed by UX3.  Thus the new pair, still
called $q$, is a genuine residual pair for $W=\varnothing$.

Project $q$ to the Levi root $-\alpha$: keep its
$r_{10},r_{32}$ entries and remove the other negative
root entries, and do the same to its local flags.
In (UX.18) this is $(x,y,t)\mapsto(x,0,0)$.
It preserves the representation, ordinary and
symplectic identities.  Replace its local $20$ flag
coordinate by the corresponding nonzero coordinate
 of $z$.  Explicitly, replace the Levi flag
 $U(m_v,0,0)$ by
 $U(m_v,0,0)U(0,z_{v,20},0)
   =U(m_v,z_{v,20},-m_vz_{v,20})$.
 The Levi flag has already diagonalized the local
 representation.  The second factor commutes with that
 diagonal matrix because \texttt{Two.atL} gives
 $\bar\chi_0=\bar\chi_2$ and symplecticity gives
 $\bar\chi_1=\bar\chi_3$ on every $G_{F_v}$.
 The resulting pair
$q'$ is in $V_0$, has all $q'_{v,20}\ne0$, and has
\texttt{SameAl} with $q$ and \texttt{SamePs} with $z$.
The $-\alpha$ representation entry is still nonzero.

The links are actual affine matrix interpolation:
\[
 p\ \xrightarrow{\ P\ }\ q\
   \xrightarrow{\ M\ }\ q'\
   \xrightarrow{\ P\ }\ z,
 \qquad A(X)=(1-X)A(0)+XA(1).
 \tag{UX.21}
\]
Here is why interpolation respects multiplication;
this cannot be assumed for arbitrary representations.
For a $P$ link, the two matrices have support $N_P$,
the same diagonal and the same $20$ entry.  The
symplectic identity also makes their $31$ entries
equal.  Their difference is supported in the
$\lambda$-off-block square-zero ideal.  Products
of any two such differences are zero, and this
ideal is preserved on both sides by the common
block triangular algebra.  For an $M$ link, both
matrices are lower triangular, have the same
diagonal and the same $10$ entry; symplecticity
also makes their $32$ entries equal.  Their
difference lies in the lower left $2\times2$
square-zero block ideal.  In either case
\[
 A(X)B(X)=(1-X)A(0)B(0)+XA(1)B(1).
 \tag{UX.22}
\]
Adjoint preserves the same ideals.  Apply (UX.22)
first to the representation identity and then to
$n(X)^{-1}r(X)n(X)$ and the symplectic identity.
It proves all pair equations over $k[X]$, the
prescribed residual diagonals and local characters,
and the stated support.  Each coefficient function
is locally constant because both endpoint functions
are.  Evaluation at $0,1$ is exactly the indicated
endpoint, so these are literal \texttt{Linked}
witnesses.

It remains to check goodness, since connectedness
alone is insufficient below.  The starting pair is
good by UX2.  If $q$ has a $-\sigma$ entry, it is
good by that entry and its nonzero $-\alpha$ entry;
otherwise its $20$ flags are nonzero whenever
needed, and it is in $V_0$ with $-\alpha$ present.
More explicitly, if the original $p$ has a
$-\sigma$ entry then \texttt{SamePs} preserves it;
if it does not, UX2 gives nonzero $20$ flags and
\texttt{SamePs} preserves those.  The pair $q'$
is good by $V_0$ and $-\alpha$.  The seed is good
by (UX.20), since \texttt{Two.bigL} makes the
$\alpha$-character nontrivial at every local group.
These are exactly the four alternatives defining
\texttt{Good}.  Thus (UX.21) proves the complete
\texttt{Chain}, uniformly for all later deep towers
and all compatible original restrictions.
\end{proof}

\subsection{Uniform bounds before choosing the deep layer}

For a residual coefficient map $a:O\to k$, write $Z(w)$ for
the locally constant scalar cocycles
\[
 z(gh)=z(g)+a(\chi_w(g))z(h)
 \tag{UX.23}
\]
which vanish on the prescribed forbidden inertia and satisfy
$z(g_0)=0$ when $w\ne0$.  Write $Z_v(w)$ for the corresponding
local cocycle space, without this normalization, and
$\operatorname{Str}(w)$ for the subspace of $Z(w)$ whose
restrictions to every $G_{F_v}$ are coboundaries.  These are
exactly \texttt{Zgl}, \texttt{Zloc} and \texttt{Str}.

\begin{proposition}[UX5: one bound for every later cyclotomic layer]
For the tame tower $T_1$ of UX4 there is an integer $C\ge1$,
chosen before the later layer or any seed coefficient field,
such that every cyclotomic layer $T\ge T_1$ satisfies
\[
 \begin{gathered}
 s\le C,\qquad
 \dim_k\operatorname{Str}(w)\le C\quad(w\in\mathcal R),\\
 \dim_k Z(w)\le C\quad(w=0,\sigma,-\sigma),\\
 \dim_k Z(w)\le d+C,\qquad
 \dim_k\prod_{v\mid\ell}Z_v(w)\le d+C
       \quad(w\in\mathcal R_o).
 \end{gathered}
 \tag{UX.24}
\]
This holds for every field $k$ and every $a:O\to k$ killing
the original maximal ideal.  One can choose the layer deep
enough that simultaneously
\[
 d\ge30(C+1)^2,\qquad d_v=[F_v:\mathbb Q_\ell]=d/s
                 \ge30(C+1).
 \tag{UX.25}
\]
Thus the literal \texttt{Bounds} used later has one
constant, independent of all later coefficient extensions.
\end{proposition}

\begin{proof}
First use the finite residue field $k_0=O/\mathfrak m_O$.
Put
\[
 H_1=G_{F_1}\cap\ker\bar r_D\cap\ker\bar c_\ell
       \ \subset G_{\mathbb Q}.
 \tag{UX.26}
\]
It is open and contains the commutator, since $F_1/\mathbb Q$
and all four residual characters are abelian.  Its fixed
field is therefore abelian and contains $\mu_\ell$.
For a layer of index $m$, the analogous subgroup is
$H_1\cap\{c_\ell\equiv1\bmod\ell^{m+1}\}$.  Indeed, on
$1+\ell\mathbb Z_\ell$, raising to $\ell-1$ is an
automorphism: in
$(u^{\ell-1}-1)/(u-1)=1+\cdots+u^{\ell-2}$ the second
factor is a unit.  This identifies the actual fixed fields
with the cyclotomic root layers, not just with abstract
fields of the same degree.  TW1 supplies a single bound $r$
for their mod-$\ell$ class ranks and hence for their
everywhere unramified additive characters.

There is also a single bound $g$ on the number of places
over $\prod_{p\in D.S}p$ in all these layers.  At $\ell$,
the local cyclotomic image is open, so the number of
places is eventually constant.  At $p\ne\ell$, the image
of Frobenius in the cyclotomic $\mathbb Z_\ell$-quotient
is open: its value is the nonzero $\ell$-adic logarithm
of $p^{\ell-1}$.  It cannot be zero, for a positive
rational prime is not an $\ell$-adic root of unity.
Multiplying these finite indices by the degree of the
fixed initial abelian field bounds the number of
decomposition-group orbits.  Sum over the finitely
many primes in $D.S$.  In particular $s\le g$.

Let $H$ be the subgroup (UX.26) in $G_F$.  Its index is
prime to $\ell$: the residual diagonal characters have
values in $k_0^\times$, and $\bar c_\ell$ has values in
$\mathbb F_\ell^\times$.  Every scalar root character
is trivial on $H$.  Restrict a strict cocycle to $H$.
It is an additive character and is unramified
everywhere.  At places over $\ell$ this follows from
the local coboundary formula, since the character is
$1$ on $H$; away from $\ell$ it follows from the
specified killed inertia.  Its coordinates in an
$\mathbb F_\ell$-basis of $k_0$ are therefore among
the $r$-dimensional space just bounded.
Restriction, together with one element of $k_0$,
determines the original cocycle.  Explicitly, if
$z|_H=0$ and $\sum_{x\in G_F/H}z(x)=0$, then
permuting the cosets by $h$ in the cocycle identity
gives $[G_F:H]z(h)=0$, hence $z(h)=0$.
Consequently
\[
 \dim_{k_0}\operatorname{Str}(w)\le r+1.
 \tag{UX.27}
\]

Here is the corresponding unrestricted global estimate.
Use DX3's finite Greenberg--Wiles equality with the full
$H^1$ local condition at every place above $D.S$.
Its dual condition is zero.  The dual strict Selmer
classes inject, after restriction to $H$, into the
everywhere unramified additive characters just bounded;
the preceding coset argument allows the same extra
factor $|k_0|$.  Local Tate duality and Euler
characteristic give, for a one-dimensional coefficient
module $M$ and a finite place $v$,
\[
 \frac{|H^1(F_v,M)|}{|H^0(F_v,M)|}
 \le |M|^{\,1+[F_v:\mathbb Q_\ell]}
       \quad(v\mid\ell),\qquad
 \frac{|H^1(F_v,M)|}{|H^0(F_v,M)|}\le |M|
       \quad(v\nmid\ell).
 \tag{UX.28}
\]
Use $|Z^1|\,|H^0|=|M|\,|H^1|$ in the global equality.
After cancelling the local $H^0$ factors and using
$|H^0(M^*(1))|\ge1$, it yields
\[
 |Z^1(G_{F,S},M)|\,
       \prod_{v\mid\infty}|H^0(F_v,M)|
       \le |M|^{r+2+g+d}.
 \tag{UX.29}
\]
This also proves finiteness; all local and global
finite-coefficient interfaces in this calculation are
the actual ones supplied by BX8.  The cocycles in $Z(w)$
inject into these $G_{F,S}$ cocycles.  For
$w=0,\pm\sigma$, complex conjugation acts trivially:
$\Psi$ is even, and the $\sigma$ root is $\Psi$ or its
inverse after relabelling.  Since $F$ is totally real,
the product in (UX.29) is then $|M|^d$.  Thus
\[
 \dim Z(w)\le r+2+g\ (w=0,\pm\sigma),\qquad
 \dim Z(w)\le d+r+2+g\ (\hbox{all }w).
 \tag{UX.30}
\]
Finally (UX.28), together with
$|Z^1|=|M||H^1|/|H^0|$, bounds each local cocycle
dimension by $2+[F_v:\mathbb Q_\ell]$.  Summing gives
$d+2s\le d+2g$.

Take $C=r+2+2g$.  These estimates prove (UX.24) over
$k_0$.  To pass to an arbitrary $k$ use its specified
embedding of $k_0$.  Every $k_0$-linear functional
$\lambda:k\to k_0$ carries each of the defining
cocycle, normalization and local-coboundary equations
to the same equation over $k_0$; in the strict case
apply $\lambda$ also to each local coboundary constant.
This preserves the bounds on dimension.  For
completeness, if $N+1$ functions over $k$ were linearly
independent, some $N+1$ evaluation vectors would have
nonzero determinant.  Expanding the finitely many
entries in a $k_0$-basis of their span and taking
coordinate functionals would put their span inside
the scalar extension of the at most $N$-dimensional
space over $k_0$, a contradiction.  The same argument
applies to functions on the disjoint union of local
groups.  No finiteness of $k$ is required.

Choose the deep layer by UX1 only now.  Every place of
the abelian field $F$ above $\ell$ has the same local
degree, so $\sum_{v\mid\ell}[F_v:\mathbb Q_\ell]=d$
gives $d_v=d/s$.  The first inequality of (UX.25)
and $1\le s\le C$ imply its second.  This order of
choices leaves (UX.24) valid after the seed and every
later coefficient enlargement.
\end{proof}

\subsection{An actual common space and its local deformation rings}

\begin{proposition}[UX6: the space, its lines, and its flat completions]
Fix a fit tower $T$ and any finite coefficient DVR
$\Lambda=O'\supset O$.  There exists the literal
$X:\mathtt{Space}(T.\mathtt{fr}\,\varnothing)(O\to\Lambda)$.
Thus $B=X.B$ is complete Noetherian local over $\Lambda$
with residue $k_\Lambda$, $A=X.A$ is a finite type
$B$-algebra, and $B\hookrightarrow A$.  The space carries
the actual pair $P$, its characteristic polynomials
$c_g(X)\in B[X]$, and its actual local characters
$\psi_{v,i}\in B$, with every clause of \texttt{Space}.
Every residual pair over $k_\Lambda[X]$ has an actual
map $A\to k_\Lambda[X]$ giving it and killing
$\mathfrak m_B$.  If $x:A\to k_\Lambda$ is a residual
pair, the completion
\[
 R_x=\widehat{A_{\ker x}},\qquad f_x:A\longrightarrow R_x
 \tag{UX.31}
\]
is flat over $A$ and is the universal complete local
ring of lifts of that exact residual pair, including
all its actual flags.
\end{proposition}

\begin{proof}
We construct all objects, since a universal pseudorepresentation
would not supply the local flags or the polynomial lines.
The residual root characters are locally constant and
have pairwise distinct values at $g_0$.  They satisfy
$\bar\chi_i\bar\chi_{i^{\rm rev}}=\bar\nu$,
$2\ne0$, and are trivial on the prescribed killed
inertia and on each $\mu_v$.  The image of $\mu_v$
in the local abelianization is $\ell$-primary torsion
and lies in abelian inertia; its image in the prime-to-$\ell$ group
$k_\Lambda^\times$ is trivial.
The quotient of $G_F$ by the closed normal subgroup
generated by the forbidden inertia is a quotient of
$G_{F,S}$ for the fixed finite set of places over
$D.S$.  It satisfies $\Phi_\ell$: additive characters
on any open subgroup inject into the finite $H^1$
over the corresponding number field, by BX8.
Each closed decomposition group satisfies the local
version by the same chapter.  These checks give
exactly \texttt{Std}, on the actual groups.

First disregard the flags and regard the $4s$ local
characters as independent.  FP1--FP2's finite-entry
construction applies to the split residual representation
$\operatorname{diag}(\bar\chi_i)$ on this quotient.
Concretely, finite $\Phi_\ell$ distinguishes its
dual-number lifts by finitely many group entries.
The finite-level entry algebras, with multiplication
and continuity imposed, have an inverse limit which
is a quotient of a power-series ring on those
tangent generators.  Impose the further closed
entry equations
\[
 \rho(g)^*\,\rho(g)=\nu(g)1,\quad
 \rho(n)=1\ (n\in N),\quad
 \rho(g_0)_{ij}=0\ (i\ne j).
 \tag{UX.32}
\]
Their ideal is closed because the power-series ring
is Noetherian and complete.  The split residual
point satisfies them, so this is a proper quotient.
Construct each independent local character by the
same argument in rank one and take the completed
coproduct of this finite list of rings.  The result
is a complete Noetherian local $\Lambda$-algebra
$R^{\rm sp}$ representing the actual matrices and
local characters with split residual matrices.
This is a direct application of the entry construction,
not an assertion that the extra hypotheses of
the full FP8 theorem hold here.

The two-dimensional diagonal symplectic torus acts
by conjugation on these matrices, fixing their
diagonal and all local characters.  It makes entry
$(i,j)$ homogeneous of weight $w_{ij}$.
Choose a finite collection of entry differences
and local-character differences spanning the
cotangent space.  They are homogeneous, vanish
at the split point, and give a surjection
$\Lambda[[z_1,\ldots,z_h]]\twoheadrightarrow R^{\rm sp}$
with homogeneous kernel; the formal torus argument
and successive lifting of homogeneous generators
are those proved in FP5.  Put
\[
 B_0=(R^{\rm sp})_0,\qquad
 A_0=\bigoplus_{w\in\mathbb Z^2}(R^{\rm sp})_w.
 \tag{UX.33}
\]
Here a homogeneous piece permits an infinite convergent
sum of monomials of its fixed weight, whereas the
direct sum in $A_0$ is a finite sum of such pieces.
The finite-monomial argument of PX1 applies here.
Every antichain in $\mathbb N^h$ is finite: from
an infinite sequence successively extract subsequences
nondecreasing in each coordinate.  Thus there are
finitely many componentwise minimal nonzero
zero-weight exponents.  They generate the whole
zero-weight monoid, by subtracting one such exponent
and inducting on total degree.  It follows that
$B_0$ is a quotient of a finite power-series ring
over $\Lambda$.  For any fixed weight, the same
antichain argument gives finitely many componentwise
minimal exponent vectors;
every other vector differs from one of them by
a zero-weight vector.  Therefore each piece is
generated over $B_0$ by finitely many monomials,
and $A_0=B_0[z_1,\ldots,z_h]$.
The degree-zero ring is complete local, with the
same residue field, and $A_0$ is finite type over it.
The same assertions descend through the homogeneous
kernel by taking its homogeneous components.

The characteristic polynomials and local characters
have degree zero.  We need their precise uniqueness
property for maps out of $B_0$.  Suppose
$u,u':B_0\to R'$ agree on these values, for a
complete local $\Lambda$-algebra $R'$ with the
specified residue field.  The diagonal eigenvalues
$d_i=\rho(g_0)_{ii}$ are determined individually by
the common polynomial $c_{g_0}$ and their distinct
residues.  Indeed
$\prod_j(u'(d_i)-u(d_j))=0$, and every factor
with $j\ne i$ is a unit; hence $u'(d_i)=u(d_i)$.
Their differences and their inverses are also
fixed by $u,u'$.

This determines every weight-zero product of
matrix entries.  Here is an explicit symplectic
invariant argument.  Let $\mathcal R$ be the algebra
generated by the matrices $\rho(g)$ over the
equalizer of $u,u'$.  It is closed under symplectic
adjoint, because $\rho(g)^*=\nu(g)\rho(g^{-1})$.
All its traces belong to that equalizer, and
\[
 e_{ii}=\prod_{j\ne i}
           \frac{\rho(g_0)-d_j1}{d_i-d_j}
                  \ \in\mathcal R .
 \tag{UX.34}
\]
In a zero-weight product, take a factor of nonzero
weight $w$.  Some other weight $w'$ has
$w\cdot w'<0$, since all weights sum to zero.
Inspection of the displayed matrix (UX.2) gives
\[
 w_{ij}\cdot w_{kl}<0\quad\Longrightarrow\quad
 j=k\ \hbox{or}\ l=i\ \hbox{or}\
 j=l^{\rm rev}\ \hbox{or}\ k^{\rm rev}=i .
 \tag{UX.35}
\]
In the first two cases multiply through the
intermediate projector $e_{jj}$ or $e_{ii}$;
in the other two first apply adjoint to one
factor.  Up to its explicit sign, the product
of the two selected entries becomes one entry
of a matrix in $\mathcal R$ of weight $w+w'$.
Induction on the number of factors reduces to
diagonal entries, which are
$\operatorname{tr}(e_{ii}M)$ and belong to the
equalizer.  This proves the assertion for every
zero-weight monomial.  The local-character
generators already belong to the equalizer.
Continuity and the finite power-series presentation
of $B_0$ now give $u=u'$.  This is the full
\texttt{genB} property; a generic-point trace
identity alone would not suffice.

We next explain why $A_0$ describes non-split
triangular residual matrices too.  If the support
is $N_P$, use the functional
$\tau(a,b)=-2a-b$; if it is $N_M$, use
$\tau(a,b)=-a-2b$.  It is strictly positive on
the permitted nonzero residual entries.  Over an
Artinian local target $R'$, write
$\mathfrak m^N=0$, choose $c$ bounding the negative
values $\tau(w_{ij})$, and form the subring
\[
 S=R'[[t]][\,\mathfrak m\,t^{-c}\,]
                         \subset R'((t)).
 \tag{UX.36}
\]
It is a finite $R'[[t]]$-module because every
product of $N$ negative generators vanishes.
It is complete Noetherian local with residue
$k_\Lambda$ and maximal ideal generated by $t$
and the nilpotent negative generators.  Rescale
entry $(i,j)$ of the given matrix by
$t^{\tau(w_{ij})}$ and leave local characters
unchanged.  A negative exponent has coefficient
in $\mathfrak m$, so all entries lie in $S$;
every off-diagonal entry lies in its maximal
ideal.  The identities
$w_{ij}+w_{jk}=w_{ik}$ and
$w_{j^{\rm rev}i^{\rm rev}}=w_{ij}$ show that
matrix multiplication and symplectic adjoint
respect this rescaling.  It therefore gives a
split residual point of $R^{\rm sp}$ over $S$.

Its value on a homogeneous element of weight $w$
is a single Laurent monomial $t^{\tau(w)}b$.
For infinite homogeneous series, this follows
by truncation: the coefficient of a fixed
Laurent power vanishes on a sufficiently high
power of the maximal ideal of $S$.
Taking the coefficient $b$ on each homogeneous
piece is multiplicative, and hence defines
$A_0\to R'$ giving the original matrices and
characters.  For a general complete local
target, do this modulo each power of its
maximal ideal.  The maps agree on the chosen
entry generators and therefore are compatible;
their limit gives the required map.  Thus
no substitution $t=1$ in an arbitrary formal
power series has been used.

Adjoin to $A_0$ the four lower unipotent
symplectic coordinates at each place and impose
the equations that the flags triangularize
the representation with exactly the chosen
local characters, together with their $\mu_v$
conditions.  These are polynomial entry
equations; the resulting algebra is finite
type over $B_0$.  Replace $B_0$ by its image
$B$ in this algebra $A$ to ensure injectivity.
The split residual pair with identity flags
gives a point, so the image quotient is
proper and is still complete local.  The
matrix and flag entries generate $A$ over
$B$, and the preceding construction gives
the required universal map for every lift
of every permitted residual pair.  Its
entries modulo $\mathfrak m_B^n$ are locally
constant: homogeneous pieces are finite
modules over the complete ring $B$, and
this is inherited from the original
continuous formal matrices.  All maps,
residual polynomials and character functions
are consequently the ones required in
\texttt{Space}.

For the line property, embed $k_\Lambda[X]$
in $k_\Lambda[[X]]$ and apply this universality
to the given polynomial residual pair.  Its
characteristic polynomials are the fixed
split polynomials and its local characters
are the fixed $\bar\chi_i$, because its
permitted support is triangular.  The just
proved uniqueness for maps from $B$ implies
that $B$ maps to constant residues.  The
generators of $A$ map to the actual
polynomial entries and flags, so the map
lands in $k_\Lambda[X]$ and kills
$\mathfrak m_B$.  Its values reproduce
the given pair exactly.

Finally $A$ is Noetherian.  Localization
followed by completion at $\ker x$ is flat,
has residue $k_\Lambda$, and is complete
Noetherian local over $\Lambda$.  An actual
lift $Q$ of $x(P)$ first gives $g:A\to R'$
by the established universality.  Equality
of the residual pairs forces
$\bar g=x$: matrix and flag entries generate
$A$, and their characteristic polynomials
and local characters determine its base.
Thus $g$ is local at $\ker x$ and extends
uniquely to $R_x$.  Conversely such a map
gives $Q$.  Uniqueness again follows from
the same entry generation and base
uniqueness.  This proves the exact
\texttt{IsDef}, flatness and residue clauses
of \texttt{cmplEx}.
\end{proof}

\begin{proposition}[UX7: the exact presentation deficit]
Suppose in addition that $T$ is deep.  For every residual
pair $x$ over $k_\Lambda$, the same completed ring $R_x$
has a presentation
\[
 R_x=\Lambda[[z_1,\ldots,z_a]]/(f_1,\ldots,f_b),
 \qquad b+2+2d\le a+2s .
 \tag{UX.37}
\]
The assertion is uniform in the residual pair, in the
space $X$, and in every permitted coefficient extension.
\end{proposition}

\begin{proof}
Write $k=k_\Lambda$, $M=\mathfrak{sp}_4(k)$, and
$\rho=\operatorname{Ad}(\bar r_x)$.  For a local
flag $\bar n_v$, let $\mathfrak b_v$ be the conjugate
of the upper Borel Lie algebra by $\bar n_v$, and
let $\mathfrak n_v$ be its complementary lower
nilpotent subspace.  These have dimensions $6$ and
$4$, while $\dim M=10$ and the diagonal torus
has dimension $2$.  For example an upper
symplectic infinitesimal matrix is
\[
 \begin{pmatrix}
 a&c&d&e\\0&b&f&d\\0&0&-b&-c\\0&0&0&-a
 \end{pmatrix},
 \tag{UX.38}
\]
which proves the asserted Borel dimension directly.

Let $\Gamma$ be the quotient of $G_F$ by the
closed normal subgroup generated by
$I_w\cap\ker\bar r_D$ at every $w\nmid\ell$.
This is exactly the quotient specified by the
pair's forbidden inertia.  The representation
$\rho$ descends to it.  To record the cohomology
being used, put
\[
 C^\bullet_{\rm ord}
 =\operatorname{Cone}\left(
 C^\bullet(\Gamma,M)\longrightarrow
       \bigoplus_{v\mid\ell}
           C^\bullet(G_{F_v},M/\mathfrak b_v)
                  \right)[-1].
 \tag{UX.39}
\]
We use continuous finite cochains.  The groups in
degrees $0,1,2$ are the source groups
$\mathtt{coneH0}$, $\mathtt{H1C}$ and $\mathtt{coneH2}$,
respectively.  A degree-one class is represented by $(z,(Y_v))$
with $z|_{G_{F_v}}-dY_v\in\mathfrak b_v$,
modulo global coboundaries and changes
$Y_v\mapsto Y_v+B_v$ with $B_v\in\mathfrak b_v$.
In degree two use $(c,(c_v))$ with
$c|_{G_{F_v}}-dc_v\in\mathfrak b_v$.
Denote their dimensions by $h^0_o,h^1_o,h^2_o$.

We first justify the arithmetic Euler inequality
on this exact quotient:
\[
             h^2_o+h^0_o+6d\le h^1_o+4d .
 \tag{UX.40}
\]
The prime-to-$\ell$ index of
$K=\ker\bar r_D|_{G_F}$ is essential here;
the quotient is not silently replaced by
the group unramified outside $\ell$.
Choose a finite set $S'$ including the places
over $\ell$ and all ramification of $K$.
Then $\Gamma$ is a quotient of $G_{F,S'}$,
by the normal closure of the indicated
inertia intersections at $S'\setminus\{\ell\}$.
At such a place, the remaining inertia has
order prime to $\ell$.  Its higher cohomology
vanishes by averaging, and the residual
unramified quotient is procyclic.  In
particular, an inflated $H^2(\Gamma,M)$ class
has zero restriction at this place.
The local tame relation and the two-term
procyclic complex give
\[
 \dim_k\operatorname{Hom}_{G_{F_w}}
       (I_w\cap K,M)
       =\dim_k H^2(F_w,M)
            \quad(w\in S',\ w\nmid\ell);
 \tag{UX.41}
\]
the homomorphisms here are locally constant,
additive, and equivariant for conjugation.
One can see (UX.41) by removing the
pro-prime-to-$\ell$ wild inertia, then using
$\operatorname{Fr}\tau\operatorname{Fr}^{-1}
=\tau^{\#k_w}$ on the pro-$\ell$ tame
generator.  The equivariant homomorphisms
are the kernel of
$\operatorname{Fr}-\#k_w$ on $M^{I_w/K}$;
local Tate duality identifies its dimension
with $H^2(F_w,M)$.  Passing through the
finite prime-to-$\ell$ inertia index changes
neither side.

The inflation five-term sequence bounds the
kernel in degree two by the sum of the spaces
in (UX.41); restriction to the inertia
generators is injective because they normally
generate the kernel.  Combining this with
global Euler characteristic and the
Poitou--Tate tail for $G_{F,S'}$ cancels
these extra-place terms.  More explicitly,
write $h^2_{\rm str}$ for the kernel of
$H^2(\Gamma,M)\to\bigoplus_{v\mid\ell}H^2(F_v,M)$.
The resulting inequality is
\[
 \begin{split}
 h^2_{\rm str}+\sum_{v\mid\ell}h^2(F_v,M)
       +h^0(\Gamma,M)+d\dim M
 \ \le\ h^1(\Gamma,M)
       +\sum_{w\mid\infty}h^0(F_w,M)
       +h^0(F,M^*(1)).
 \end{split}
 \tag{UX.42}
\]
For clarity about the cancellation, the five-term
sequence gives
$\dim\ker(\mathrm{inf}:H^2(\Gamma,M)\to H^2(G_{F,S'},M))
=\dim H^1(\ker,\!M)^{\Gamma}
 -h^1(G_{F,S'},M)+h^1(\Gamma,M)$.
Its first term is at most the sum in (UX.41).
The strict kernel maps into the strict kernel
for all places of $S'$, by the local vanishing
above.  Apply the Poitou--Tate tail to that
kernel and insert the global Euler equality;
the same sum of $h^2(F_w,M)$ is subtracted
on the other side.  This proves (UX.42)
and finiteness, with no omitted local term.

The last term of (UX.42) vanishes.  The
triangular support of the residual pair
gives a filtration of the matrix adjoint
module with characters $\bar\chi_i/\bar\chi_j$.
No one of these is $\bar c_\ell$, by
$T.\mathtt{Ok.g2}$.  A homomorphism to
$k(1)$ must therefore vanish, successively
on the filtration.  This also applies
to $M$, since
$X\mapsto(X-X^*)/2$ is an equivariant
projection from all matrices onto $M$.
At a real place, the involution $\bar r_x(c)$
has multiplier $-1$; its $+1$ and $-1$
eigenspaces are paired perfectly by the
symplectic form and both have dimension $2$.
The fixed part of $M$ consists of an
endomorphism on the first eigenspace and
the forced negative transpose on the second,
and has dimension $4$.

Apply now the exact sequence of (UX.39).
Local Tate duality makes
$H^2(F_v,M)\to H^2(F_v,M/\mathfrak b_v)$
surjective, by dualizing the injective
map on $H^0$ of the corresponding duals.
Elementary kernel counting gives
\[
 \begin{split}
 &\dim\ker\left(H^2(\Gamma,M)\to
             \bigoplus_vH^2(F_v,M/\mathfrak b_v)\right)
            +\sum_vh^2(F_v,M/\mathfrak b_v)\\
 &\hspace{25mm}\le
 h^2_{\rm str}+\sum_vh^2(F_v,M).
 \end{split}
 \tag{UX.43}
\]
Local Euler characteristic gives
$h^1(F_v,Q_v)-h^0(F_v,Q_v)
 =h^2(F_v,Q_v)+d_v\dim Q_v$
for $Q_v=M/\mathfrak b_v$.  Subtract consecutive
ranks in the cone exact sequence and use
(UX.42)--(UX.43), $\dim Q_v=4$,
$\sum d_v=d$, and the vanishing dual
invariants.  The result is (UX.40).

Temporarily remove the $\mu_v$ conditions
from the deformation problem and call its
universal completed ring $R^0$.  Its
tangent dimension $a$ satisfies
\[
                     h^1_o+2\le a+h^0_o.
 \tag{UX.44}
\]
Indeed every cone class has a representative
whose local $Y_v$ lies in the complement
$\mathfrak n_v$.  Conjugation by $1+\epsilon X$
can make the off-diagonal value of $z(g_0)$
zero, since the residual eigenvalue
differences are units.  This is the actual
normalization imposed in (UX.32).
The remaining global conjugations are the
two diagonal directions; those acting
trivially are exactly the invariants lying
in every $\mathfrak b_v$, of dimension $h^0_o$.
This proves (UX.44), and in these separated
coordinates it in fact gives equality.

The space $H^2_o$ is a complete obstruction
space for $R^0$.  Here is its concrete
realization.  Across a small extension
with square-zero kernel $I$, lift the
global matrices and the local flags
preserving their symplectic form.  Such
lifts exist: the correction
$X\mapsto(X-X^*)/2$ solves the infinitesimal
form equation, and the lower unipotent
chart lifts its four coordinates freely.
The multiplication defects of the global
matrices form a continuous two-cocycle
in $M\otimes I$.  The failure of each
lifted flag to triangularize the local
matrices gives the local one-cochain
modulo $\mathfrak b_v\otimes I$.
Associativity and restriction give exactly
the cone equations above.  Changing
lifts changes this pair by a cone
coboundary.  If its class is zero,
correct the matrices by the corresponding
global one-cochain and correct the
flags by its local zero-cochains.
The remaining $g_0$ normalization is
restored by the separated-eigenvalue
conjugation, which preserves all other
conditions.  Thus vanishing is sufficient
as well as necessary for a lift.

The finite obstruction argument of FP7
now presents $R^0$ with $a$ variables
and at most $h^2_o$ relations.  Explicitly,
in a minimal surjection $\Lambda[[z_a]]
\twoheadrightarrow R^0$, dualizing the
relation module $J/\mathfrak mJ$ sends a
functional to the obstruction for the
associated one-dimensional small
extension.  If that obstruction vanishes,
the corrected universal lift gives a
section, so the functional on $J$ is zero.
This injects the dual relation module
into $H^2_o$ and proves the bound.

Finally impose $\mu_v$.  By local class
field theory, its image in the abelianization
of $G_{F_v}$ is exactly the $\ell$-power
torsion in the profinite completion of
$F_v^\times$, which is cyclic: it is
the group of $\ell$-power roots of unity.
Choose a lift $\zeta_v$ of a generator.
All characters kill the closed commutator,
so it suffices to impose their values
at $\zeta_v$.  Deepness gives
$\nu(\zeta_v)=1$: the cyclotomic character
has no $\ell$-power torsion for odd $\ell$,
the dyadic character has prime-to-$\ell$
order, and the remaining $\psi|_{G_F}$
does also.  The symplectic relation
$\psi_{v,i}\psi_{v,i^{\rm rev}}=\nu$
therefore reduces the four equations to
\[
             \psi_{v,0}(\zeta_v)=1,\qquad
             \psi_{v,1}(\zeta_v)=1 .
 \tag{UX.45}
\]
Thus $R_x$ is a quotient of $R^0$ by
at most $2s$ further equations.  If
$b\le h^2_o+2s$ is the resulting count,
(UX.40) and (UX.44) give
$b+2+2d\le a+2s$, as required.
\end{proof}

\subsection{The two loci that a connecting curve must avoid}

For reference, the nine source reducible patterns are the following.
The notation in each row lists all its permitted nonzero roots.
\[
\begin{array}{c|l}
0&\beta,-\beta,-\alpha,-\tau,-\sigma\\
1&\beta,-\beta,\alpha,\tau,\sigma\\
2&\tau,-\tau,\sigma,\beta,-\alpha\\
3&\tau,-\tau,-\sigma,-\beta,\alpha\\
4&\alpha,-\alpha,-\tau,-\sigma,-\beta\\
5&\alpha,-\alpha,\tau,\sigma,\beta\\
6&\sigma,-\sigma,\beta,-\alpha,-\tau\\
7&\sigma,-\sigma,-\beta,\alpha,\tau\\
8&\beta,-\beta,\tau,-\tau
\end{array}
\tag{UX.46}
\]
Write $I_S$ for the ideal generated by the representation
entries outside row $S$.  Over a field with separated
$g_0$, every non-absolutely-irreducible representation
lies in one of these loci.  To check the assertion,
the spectral projectors (UX.34) lie in the algebra
spanned by the representation.  Multiplying between
them shows that this algebra is spanned by its
nonzero matrix units.  If their directed graph is
strongly connected, multiplying along paths gives
all matrix units and the representation is absolutely
irreducible.  Otherwise a union of terminal components
gives a proper coordinate invariant subspace.
The algebra is closed under symplectic adjoint,
so an edge $(i,j)$ is present if and only if
$(j^{\rm rev},i^{\rm rev})$ is present.  List the
nonempty proper coordinate subsets and their
symplectic orthogonals.  The four one-dimensional
subsets give rows $0$--$3$, the four isotropic
two-dimensional subsets give rows $4$--$7$,
and the two complementary nonisotropic subsets
give the common block pattern in row $8$.
Three-dimensional subsets give the orthogonals
of the one-dimensional ones.  This exhausts
the list and proves the assertion without a
classification assumption.

\begin{proposition}[UX8: the complete reducible-locus bound]
At every good residual pair $x$ and for every row $S$
of (UX.46), the exact ring of UX6 satisfies
\[
                   \dim(R_x/I_S)\le d+20C+1 .
 \tag{UX.47}
\]
Here $C$ is the one constant of UX5; no bound on the
cardinality of the residue field is added.
\end{proposition}

\begin{proof}
There is nothing to prove if $I_S=R_x$.  Otherwise
the residual pair $c=(\bar r,\bar n)$ itself has
support in $S$.  A derivation at its residual
point gives the infinitesimal matrices
\[
 X(g)=\partial r(g)\bar r(g)^{-1},\qquad
 Y_v=\partial n_v\,\bar n_v^{-1}.
 \tag{UX.48}
\]
They lie in $\mathfrak{sp}_4(k)$, satisfy the
adjoint cocycle identity and vanish on the
forbidden inertia, and $X(g_0)$ is diagonal.
They also satisfy
\[
 \begin{gathered}
 (X(g)\bar r(g))_{ij}=0\quad
                       (i\ne j,\ w_{ij}\notin S),\\
 \bar n_v^{-1}Y_v\bar n_v\text{ is strictly lower},\\
 \bar n_v^{-1}\!
 \left(X(h)+\bar r(h)Y_v\bar r(h)^{-1}-Y_v\right)
 \bar n_v\text{ is upper}\quad(h\in G_{F_v}).
 \end{gathered}
 \tag{UX.49}
\]
Conversely these are exactly the linearized
equations needed for an upper bound.  The
map from derivations to $(X,Y)$ is injective:
if all entries in (UX.48) vanish, the two
dual-number maps give identical pairs;
UX6's universal uniqueness makes the maps
equal.  Thus it is enough to bound the
space defined by (UX.49).

We reduce this linear space to two split
calculations.  Conjugate matrices by
$\operatorname{diag}(t^{m_i})$ and take
the limit at $t=0$, using respectively
\[
 (m_i)=(0,2,1,3)\text{ for }N_P,\qquad
 (0,1,2,3)\text{ for }N_M.
 \tag{UX.50}
\]
The residual representation becomes split
and the flags become $1$.  If goodness
instead uses $V_0$ and its nonzero $20$
flag coordinates, use $(m_i)=(0,1,0,1)$.
The representation again becomes split,
but the flags become
$1+b_vX_{-\sigma}$, with $b_v=c.n_v{}_{20}\ne0$.
All these weights satisfy
$m_i+m_{i^{\rm rev}}$ constant, so they
preserve the symplectic equations.
The dimension of the original tangent
space is at most that of the limiting
space.  This assertion does not require
a finite-dimensional ambient space of
all functions: filter the function
coordinates by their finitely many
matrix weights.  The lowest nonzero
weight terms of the linear equations
(UX.49) give the limiting equations.
Each successive graded piece injects
into the corresponding weight piece
of their solution space.  Summing the
finitely many inequalities gives the
assertion.

For split $\bar r$, every scalar
$X_{ij}$ belongs to $Z(w_{ij})$.
When $\bar n_v=1$, every lower entry
is locally a coboundary, directly by
the last equation in (UX.49).
If the first or second goodness
alternative holds, only rows $0,4,6$
can contain its required entries.
In these three rows select, respectively,
the entries $12,01,12$.  Upon setting
the selected scalar cocycle equal to
zero, every other independent odd
entry is either lower or is forced
to be zero by the pattern.  Thus all
odd entries lie in their strict
spaces $\operatorname{Str}(w)$.

In the $V_0$ case, the local diagonal
characters satisfy $\bar\chi_0=\bar\chi_2$
and $\bar\chi_1=\bar\chi_3$.  The first
equality follows already from the $20$
entry of $\bar r(h)\bar n_v=
\bar n_v(\bar n_v^{-1}\bar r(h)\bar n_v)$,
using $b_v\ne0$; the second follows by
the multiplier.  In particular, on
this local group a root's character
depends only on its first coordinate.
Set
$Z=X(h)+\bar r(h)Y_v\bar r(h)^{-1}-Y_v$.
Multiplying out the limiting flag
$1+b_vX_{-\sigma}$ in (UX.49) gives
the three relations
\[
 Z_{10}=-b_vZ_{12},\qquad
 Z_{30}=-b_v^2Z_{12},\qquad
 Z_{21}=2b_vZ_{01}+b_v^2Z_{03}.
 \tag{UX.51}
\]
The possible rows are $0,2,4,5,6,8$;
choose the single exceptional entry as follows:
\[
 \begin{array}{c|rrrrrr}
 S&0&2&4&5&6&8\\ \hline
 (i,j)&12&12&01&01&12&03 .
 \end{array}
 \tag{UX.52}
\]
If the selected cocycle is zero, (UX.51)
and the entries excluded by that row
make all six independent odd entries
locally coboundaries.  For example the
first relation transfers this property
between $X_{12}$ and $X_{10}$, and
the second transfers it between
$X_{12}$ and $X_{30}$.  In the third,
any two known entries determine the
third.  Division by $b_v,b_v^2,2b_v$
is permitted.  These statements concern
coboundaries, not zero values of $Z$:
the same identities absorb the local
$Y_v$ entries into the required
coboundary constants.  This verifies
every row of (UX.52).

Project the whole tangent space to the
selected odd scalar cocycle.  Its image
has dimension at most $d+C$, by UX5.
In the kernel use the ten entries
$i\le j^{\rm rev}$ of a symplectic
infinitesimal matrix.  Each odd one
is strict, and each even one has
weight $0,\sigma$ or $-\sigma$;
each therefore contributes at most
$C$.  The local variables contribute
at most their six strictly lower
coordinates at each place, hence
at most $6s\le6C$.
These coordinates determine $(X,Y)$:
symplectic adjoint determines the
remaining matrix entries, and
conjugating back by $\bar n_v$
determines each $Y_v$.
Consequently
\[
 \dim\operatorname{Tan}_S(c)
        \le(d+C)+10C+6C
        \le d+20C.
 \tag{UX.53}
\]
The first or second goodness alternatives
use (UX.50); the third or fourth use
(UX.52).  Thus all good residual pairs
are covered.  A complete local algebra
over a DVR with relative cotangent
dimension at most $d+20C$ is a quotient
of a power-series ring in that many
variables over the DVR.  Its dimension
is therefore at most $d+20C+1$.
Apply this to $R_x/I_S$ and the
injectivity in (UX.48), proving (UX.47).
\end{proof}

\begin{proposition}[UX9: trace fibres and coincident local characters]
Let $J\subset B$ satisfy the literal five-field
promodular-support contract for $T,X,\iota$.
For a prime $\mathfrak q\subset R_x$ at which the
generic representation is absolutely irreducible,
put $\mathfrak q_B=\mathfrak q\cap B$.  Then
\[
       \dim R_x/\mathfrak q
             \le \dim B/\mathfrak q_B+2+4s .
 \tag{UX.54}
\]
If a prime $\mathfrak b\subset B$ contains
$J+\mathfrak m_\Lambda B$ and, for one $v$ and
$i\ne j$, all
$\psi_{v,i}(h)-\psi_{v,j}(h)$ with $h\in T.U_1(v)$,
then
\[
                    \dim B/\mathfrak b+d_v\le2d+1 .
 \tag{UX.55}
\]
Only the \texttt{fin} field of the support contract
is used in (UX.55).
\end{proposition}

\begin{proof}
First let $L$ be the fraction field of the image
of $A$ in $R_x/\mathfrak q$, and let $E$ be
the fraction field of $B/\mathfrak q_B$.
All representation traces and multipliers lie
in $E$.  Absolute irreducibility implies that
the algebra spanned by its matrices is $M_4(L)$.
Choose sixteen matrices forming a basis.  The
trace pairing on $M_4$ is nondegenerate, so
the inverse of their Gram matrix expresses
all multiplication structure constants and
all other representation matrices using
traces of products.  They lie in $E$.
Thus the representation algebra descends to
a central simple $E$-algebra of dimension $16$.
After a finite algebraic extension of $E$
it splits, and the representation matrices
are simultaneously defined over that
extension up to conjugation.

The unique alternating form with the
specified multiplier is defined there up
to scalar: solve its linear entry equations
over that field and use irreducibility
to see that their solution space is one
dimensional.  A symplectic basis and then
an ordered eigenbasis for $r(g_0)$ may be
chosen after another finite algebraic
extension.  The separating eigenvalues
are roots of $c_{g_0}\in E[X]$.  Both
the original and the new $g_0$ matrices
are diagonal with this same ordering,
so their conjugator is diagonal.
The fixed form forces
$h_0h_3=h_1h_2$ up to a scalar which
is algebraic over $E$ and can be
absorbed after adjoining its square
root.  Conjugation is insensitive to
a common scalar.  Thus only two
diagonal parameters are needed to
recover all original representation
entries over an algebraic extension
of $E$.  This argument can be performed
in an algebraic closure of $L$, which
does not change transcendence degrees.

A lower unipotent symplectic flag is
determined by its entries $10,20,30,21$:
its other nonconstant entries are
$n_{32}=-n_{10}$ and
$n_{31}=n_{20}-n_{21}n_{10}$.
Consequently
\[
                 \operatorname{trdeg}_E L\le2+4s .
 \tag{UX.56}
\]
Every entry of $A$ is generated over $B$
by these actual matrices and flags,
which is why (UX.56) bounds this field,
not just the field of traces.
The complete local ring $B$ is excellent
and universally catenary.  The dimension
formula for a finite type algebra at a
closed point with finite residue extension
gives
$\dim(A/\mathfrak q_A)_{\ker x}
 \le\dim B/\mathfrak q_B+
                 \operatorname{trdeg}_E L$.
Its completion is
$R_x/\mathfrak q_A R_x$ and has the
same dimension; its quotient
$R_x/\mathfrak q$ has no larger dimension.
This proves (UX.54), including the
completion step.

For (UX.55) put $\overline B=B/\mathfrak b$,
a complete local domain of characteristic
$\ell$.  The actual local characters
take inertia into $1+\mathfrak m_{\overline B}$
because $T.\mathtt{Ok.ineL}$ makes the
original residual representation trivial
on inertia above $\ell$.  Their products
in symplectically paired positions equal
$1$ there: the multiplier reduces to
its fixed coefficient residue.  The
torsion part of a pro-$\ell$ character
is zero over this domain, since
$u^{\ell^a}=1$ implies $(u-1)^{\ell^a}=0$.
Local reciprocity and
$\mathcal O_{F_v}^{\times}\{\text{pro-}\ell\}
 \simeq(\text{finite torsion})\times\mathbb Z_\ell^{d_v}$
therefore give $d_v$ inertia elements
which determine each such character.
The same fixed number $d_v=d/s$ works
at every place.

Also choose one inertia element $h_0$
whose cyclotomic value generates the
pro-$\ell$ part of the cyclotomic image.
Two characters into $1+\mathfrak m_{\overline B}$
which agree on $\ker c_\ell$ and on
$h_0$ agree on all inertia: their
ratio descends to a procyclic quotient,
and its prime-to-$\ell$ part is
automatically trivial.  Here
$I_v\cap\ker c_\ell=T.U_1(v)$,
since the residual kernel condition
is automatic on $I_v$.

Choose $a\in\{0,1,2,3\}$ so that
$\{a,a^{\rm rev},i,i^{\rm rev}\}$
is the whole set and either $a=j$
or $j=i^{\rm rev}$.  In $B/\mathfrak b$
use the following parameters:
the $d_v$ pinning values
$\psi_{u,a}(h)-1$ at each place $u$;
the $d_v$ values $\psi_{u,i}(h)-1$
at all $u\ne v$; and the single
value $\psi_{v,i}(h_0)-\psi_{v,j}(h_0)$.
Their number is
\[
               sd_v+(s-1)d_v+1=2d-d_v+1 .
 \tag{UX.57}
\]
They lie in the maximal ideal.  In
any prime quotient killing them,
all $\psi_{u,a}$ are $1$ on inertia,
and all $\psi_{u,i}$ with $u\ne v$
are $1$.  At $v$ the hypothesis on
$U_1(v)$ and the additional parameter
make $\psi_{v,i}=\psi_{v,j}$ on
inertia.  If $a=j$, this makes it
$1$; if $j=i^{\rm rev}$, the paired
product gives $\psi_{v,i}^2=1$,
and its residue $1$ with $2$ invertible
again makes it $1$.  The paired
characters now make all four local
characters trivial on all inertia.

The precise \texttt{fin} contract says
that some power of $\mathfrak m_B$
is contained in
\[
 J+\mathfrak m_\Lambda B+
   (\psi_{u,m}(h)-1:\ u,\ m,\ h\in I_u).
 \tag{UX.58}
\]
It therefore forces every such prime
to be the maximal ideal.  The parameters
in (UX.57) have maximal radical in
$B/\mathfrak b$.  The height theorem
bounds its dimension by their number,
proving (UX.55).  This argument neither
requires $B/J$ finite over $\Lambda$
nor replaces (UX.58) by mere generic
finiteness.
\end{proof}

\subsection{Reaching every component at a good point}

For a prime $\mathfrak p\subset A$, write
$\operatorname{SomeIn}(J,\mathfrak p)$ if some
minimal prime $\mathfrak q\subset\mathfrak p$
contains $JA$, and write
$\operatorname{AllIn}(J,\mathfrak p)$ if all
such minimal primes do.  For a map $y:A\to D$
to a domain, $\operatorname{LocZero}(J,y)$
means that every prime $\mathfrak b\subset B$
annihilated by $y|_B$ contains $J$.  These are
the exact source predicates, with their
quantifiers over primes retained.

\begin{proposition}[UX10: quantitative component propagation and residual lines]
Assume (UX.24)--(UX.25) and the literal support
contract for $T,X,J,\iota$.
At any good residual point $x:A\to k_\Lambda$,
\[
 \operatorname{SomeIn}(J,\ker x)
       \ \Longrightarrow\
 \operatorname{AllIn}(J,\ker x).
 \tag{UX.59}
\]
Moreover, if two residual points are the ends
of a finite chain of literal polynomial links
whose endpoints are good, support at one
endpoint in the sense \texttt{SomeIn} implies
\texttt{AllIn} at the other.  Only the
\texttt{fin} and \texttt{p1} fields of the
support contract are used in this proposition.
\end{proposition}

\begin{proof}
Put
\[
                    N=2d+1-2s .
 \tag{UX.60}
\]
It is a positive integer.  The numerical
choices imply the two strict inequalities
\[
 d+20C+1<N,\qquad
 2d+1+(2+4s)<N+d_v .
 \tag{UX.61}
\]
Indeed $s\le C$ reduces the first to
$d>22C$, and the second to $d_v>2+6s$;
both follow from (UX.25).

Complete at $x$, obtaining $R=R_x$.
By UX7 it is a quotient of the regular
complete local ring
$\Lambda[[z_1,\ldots,z_a]]$ by $b$
elements with $a-b\ge N+1$.
Grothendieck's connectedness theorem gives
a chain between any two minimal primes
of $R$ such that adjacent primes $P,Q$
satisfy
\[
                       \dim R/(P+Q)\ge N+1 .
 \tag{UX.62}
\]
Here is the exact specialization of the
classical theorem.  The power-series
domain has dimension $a+1$ and cannot
be disconnected by removing any proper
closed subset.  The local cohomological
dimension of a $b$-generated ideal is
at most $b$, by its \v Cech complex.
Stacks Project, Lemma 52.14.1
(\href{https://stacks.math.columbia.edu/tag/0ECR}{tag 0ECR}),
therefore says that its quotient remains
connected after removing a closed subset
of dimension less than $a-b$.
If the graph of minimal components
joined by intersections of dimension
at least $a-b$ were disconnected,
remove the finitely many cross
intersections between its two parts.
The resulting space would be disconnected,
contradicting the theorem.  This proves
(UX.62), without assuming that $R$
is a complete intersection.

We show that $JR\subset P$ propagates
across one such edge to $JR\subset Q$.
Let $\varpi$ be a uniformizer of
$\Lambda$.  The first family of bad
ideals is the nine $I_S$ of UX8.
The second consists, for every $v,i\ne j$,
of
\[
 B_{vij}=JR+(\varpi)+
     (\psi_{v,i}(h)-\psi_{v,j}(h):
                                      h\in U_1(v))R .
 \tag{UX.63}
\]
For every prime $\mathfrak q$ containing
one of these latter ideals but none of
the $I_S$, the generic representation
is absolutely irreducible by (UX.46).
UX9 and (UX.61) give
\[
        \dim R/\mathfrak q
        \le2d+1-d_v+(2+4s)<N .
 \tag{UX.64}
\]
Every quotient by an $I_S$ has
dimension less than $N$ by (UX.47)
and (UX.61).

Cut $R/(P+Q)$ by $\varpi$.
The principal ideal theorem leaves
a component of dimension at least $N$.
Choose its prime $\eta$.  It cannot
contain an $I_S$, by the preceding
bound.  Nor can it contain a $B_{vij}$,
by (UX.64).  For each bad ideal choose
one element outside $\eta$ and take
their product $f$, which remains
nonzero in the complete local domain
$R/\eta$.  There exists a
one-dimensional prime quotient of
this domain on which $f$ remains
nonzero.  For an explicit construction
in dimension $n>1$, extend $f$ to
a system of parameters
$f,t_2,\ldots,t_n$.  A minimal prime
over $(t_2,\ldots,t_n)$ avoiding $f$
exists by the height theorem; its
quotient has dimension one because
adding $f$ makes the ideal maximal.
For $n=1$ take the zero prime.
Let $p\subset R$ be its inverse
image.  Then
\[
 P+Q+(\varpi)\subset p,\qquad
 \dim R/p=1,\qquad
 I_S\not\subset p,\ B_{vij}\not\subset p .
 \tag{UX.65}
\]

The normalization of the complete
one-dimensional domain $R/p$ is
finite, since complete local rings
are Nagata.  It is local: a finite
algebra over a complete local ring
is a product of its local factors,
and a subring of a field has no
nontrivial idempotents.  Hence this
normalization is a complete DVR.
Its residue field $k'$ is a finite
extension of the finite field $k_\Lambda$.
It has characteristic $\ell$ because
$\varpi\in p$.  The equal-characteristic
complete DVR structure theorem
identifies it with $k'[[t]]$.
The coefficient-field embedding can
be chosen to extend the given
$k_\Lambda$: its finite field
elements are the unique lifts of
their simple equations
$X^{|k_\Lambda|}-X=0$.
Thus there is a genuine
$\Lambda$-algebra map
\[
 R\longrightarrow k'[[t]]
                  \quad\hbox{with kernel }p .
 \tag{UX.66}
\]
The normalization and structure
theorems used here are, precisely,
Stacks Project, Section 33.41
(\href{https://stacks.math.columbia.edu/tag/0C44}{tag 0C44})
and Lemma 10.160.10
(\href{https://stacks.math.columbia.edu/tag/0C0S}{tag 0C0S}).

Compose (UX.66) with $A\to R$.
This is one of the literal finite-field
power-series points in \texttt{p1}.
Its base maximal ideal has zero
constant coefficient, since both
maps are local.  It kills $J$ because
$JR\subset P\subset p$.
Its generic representation is
absolutely irreducible because
it avoids all $I_S$.  Every pair
of local characters is distinct
on some $U_1(v)$: otherwise all
their differences would lie in
$p$, contradicting (UX.65).
Apply \texttt{p1} to obtain
$\operatorname{LocZero}$ for this
actual point.  The prime
$Q\cap B$ is annihilated by it,
as $Q\subset p$; therefore
$J\subset Q\cap B$, as required.
No use of purity or characteristic-zero
automorphy was made in this step.

Faithful flatness of
$A_{\ker x}\to R$ lifts each
minimal prime of $A$ below
$\ker x$ to a minimal prime of
$R$, with that contraction.
To see the minimal assertion,
first use lying over for a
faithfully flat map and then
choose a minimal prime below
the lifted prime; going down
forces its contraction to be
the original minimal prime.
Support on one original
component therefore gives
support on one component
of $R$.  Propagation along
(UX.62) gives support on
every component of $R$,
and contraction gives
(UX.59).

Finally a polynomial link gives,
by UX6, a map $A\to k_\Lambda[X]$
whose evaluations are the two
actual endpoint maps.  Its
kernel is prime and is contained
in both endpoint maximal ideals.
Choose a minimal prime below
this common kernel.  If
\texttt{AllIn} holds at the
first endpoint, that minimal
prime contains $JA$, so
\texttt{SomeIn} holds at the
second.  Apply (UX.59) there.
Induction along a finite chain,
in either direction, proves
the final assertion.  This
uses goodness only at the
endpoints; it does not
assume all points of a
link are good or automorphic.
\end{proof}

\subsection{Returning to the original valued representation}

\begin{proposition}[UX11: integral restriction and exact coefficient descent]
Let the original global pair have the lattice shape of UX2,
and let $T$ be an extension of its initial tower as in UX1.
There is a small pair $P$ on $G_F$ whose representation is
the restriction of that same global representation,
whose residue is the prescribed residual restriction,
and whose local inertia characters satisfy
\[
 \psi_{v,i}(s)=a_O\bigl(D.\mathrm{ps}_i(s|_{\mathbb Q})\bigr)
 \qquad(s\in I_{F_v}).
 \tag{UX.67}
\]
For every small pair $P$ there is a coefficient threshold
$O_3$ such that, for every coefficient DVR
$\Lambda\supset O_3$, it descends along an injective map
$f:R\to V$, where $R$ is a complete Noetherian local
$\Lambda$-algebra with residue field exactly $k_\Lambda$,
and
\[
 f\circ a_{\Lambda,R}=a_\Lambda,
 \qquad \operatorname{res}_V\circ f
       =j_\Lambda\circ\operatorname{res}_R.
 \tag{UX.68}
\]
The descended pair is continuous for the maximal-ideal
topology of $R$.  Independently, every finite residual
chain in the ambient residue field descends, with all
its actual polynomial links and goodness conditions,
to $k_\Lambda$ for every $\Lambda$ above a single
coefficient threshold.
\end{proposition}

\begin{proof}
First choose a complete DVR $A$ and an injective local
map $A\to V$ witnessing smallness of the global pair.
For each of the finitely many places $v\mid\ell$, choose
the prescribed Galois element carrying the original
place over $\ell$ to the place selected in $T$.
Its representation matrix and the original ordinary flag
are matrices over $A$.  Their product has all leading
principal minors invertible.  This can be checked in
the residue field: support $N_P$ and flag support $Q$
make these minors, successively, the products of the
first one, two, three and four diagonal characters.
For instance the first three rows of their product have
the form
\[
 \begin{pmatrix}
  a&0&0&0\\ *&b&*&0\\ *&0&c&0
 \end{pmatrix},\qquad abc\ne0.
 \tag{UX.69}
\]
Gaussian elimination over $A$ therefore gives a unique
lower unipotent factor $n_v$ and an invertible upper
factor.  Both have the required symplectic similitude
properties: apply the adjoint relation to the
factorization and use uniqueness of lower-times-upper
factorization.  The lower factor has multiplier $1$.
This construction uses divisions by units only.

Conjugating the original local representation by this
factor gives an upper triangular representation.  Its
ordered diagonal characters are unchanged by the
remaining upper triangular conjugation.  Since the
original characters are abelian, conjugating the
decomposition group changes no character value.
This proves (UX.67) with the original indices $i$.
The tower label permutes residual global characters;
it does \emph{not} permute the integral inertia characters
in (UX.67).  The source's label condition is precisely
what makes the two residual descriptions agree on
these decomposition groups.

The same elimination in the residue field gives flag
support $Q$ and the residual restriction appearing in
UX4.  Restriction preserves normalization at $g_0$,
because the later $g_0$ is the specified $\ell$-power
of the earlier one.  The original unramified conditions,
together with the tower condition killing $\psi$ on
the remaining forbidden inertia, give the exact
\texttt{IsPair} inertia condition outside $\ell$.
For $\mu_v$, the image in the local abelianization is
$\ell$-primary torsion.  Local class field theory places
it in inertia; the tower's cyclotomic and character
conditions make the original inertia characters
trivial there.  Thus the required $\mu_v$ condition
also holds.  The finitely many matrices $n_v$ were
constructed in $A$, and the restricted representation
is continuous there, so the new pair is still small.

We prove the coefficient assertion without assuming that
the ambient valued field has characteristic zero.
Write $f_A:A\to V$ for the small DVR map.  The ambient
field $K$ is algebraic over $\operatorname{Frac}A$ by
the definition of smallness.  A useful consequence is
\[
 \bigcap_{n\ge0}f_A(\pi)^nV=0,
 \tag{UX.70}
\]
where $\pi$ is a uniformizer of $A$.
Indeed, if nonzero $x$ belonged to the intersection,
an algebraic relation for $x$, after removing a power
of its variable and clearing denominators, would have
nonzero constant term $u\pi^e$ with $u\in A^\times$.
Writing $x=f_A(\pi)^{e+1}z$ in that relation would
put the unit $f_A(u)$ in $f_A(\pi)V$, a contradiction.
This also explains why no unspecified continuity of
the ambient valuation is being used.

For $t\in\mathbb Z_\ell$, choose ordinary integers
$t_n\equiv t\pmod{\ell^n}$.  Their images in $A$
form a Cauchy sequence, since $\ell\in\mathfrak m_A$;
let its limit be $t_A$.  In characteristic $\ell$
the sequence is eventually constant.  Compatibility
of the given coefficient maps with integers shows
that
$f_A(t_A)-a_O(t)$ belongs to every
$f_A(\pi)^nV$.  By (UX.70) it is zero.  Consequently
the image of $\mathbb Z_\ell$ under every given
coefficient map already lies in $f_A(A)$.

The residue field of $A$ is finite.  Choose one lift
of each of its elements.  The ambient residue
approximation property puts each of their residues
in the image of some coefficient residue field.
A finite compositum gives $O_3$ such that, for every
$\Lambda\supset O_3$, all these residues belong to
$j_\Lambda(k_\Lambda)$.
Choose finitely many $\mathbb Z_\ell$-module generators
of $\Lambda$.  Adjoin their prescribed images in $K$
to $\operatorname{Frac}A$, obtaining a finite field
extension $E$.  Then
\[
 A'=E\cap V
 \tag{UX.71}
\]
is a complete DVR, finite over $A$, with finite residue
field; it contains the image of $\Lambda$ and of $A$.
Here is an elementary finiteness argument suited to
the possibly positive-characteristic ambient field.
The ring $A'$ is a valuation ring of finite $A$-rank
at most $[E:\operatorname{Frac}A]$.  Independent
residue classes in $A'/\pi A'$ lift to $A$-linearly
independent elements: divide a putative relation by
the smallest valuation of a nonzero coefficient and
reduce modulo $\pi$.  Thus $A'/\pi A'$ is finite
over $A/\pi A$.  Lift a finite spanning family and
successively express each element modulo
$\pi,\pi^2,\ldots$.  Completeness of $A$ sums the
coefficient series; (UX.70) says the remaining error
is zero.  The family therefore generates $A'$ over
$A$.  Hence $A'$ is Noetherian; a nonfield
Noetherian valuation ring is a DVR.  Finiteness
over complete $A$ gives completeness, and its residue
field is finite.  This proves all claims in (UX.71).

Let $k'$ be the residue field of $A'$, and let
$k_\Lambda\hookrightarrow k'$ be the injection
induced by the actual coefficient map.  Define
\[
 R=\{x\in A':\bar x\in k_\Lambda\}.
 \tag{UX.72}
\]
Its maximal ideal is $\mathfrak m_{A'}$ as a set,
and its $n$th power is $\mathfrak m_{A'}^n$.
For the latter assertion write any element
$\pi'^na$ as the product of $n-1$ copies of
$\pi'\in R$ and $\pi'a\in R$.
The ring $R$ is closed and complete, has residue
field $k_\Lambda$, and has finite-dimensional
cotangent space: $k'$ is finite over $k_\Lambda$.
Lift a finite spanning family of that cotangent
space.  Successive approximation, with the residue
field supplied by $\Lambda$, gives a surjection
from a power-series ring in finitely many variables
over $\Lambda$ to $R$.  Thus $R$ is Noetherian.
The choice of $O_3$ puts all of $A$ inside $R$.
Its inclusion is local, hence carries
$\mathfrak m_A^n$ into $\mathfrak m_R^n$.
The original continuous pair therefore descends
continuously to $R$.  The inclusion $R\hookrightarrow V$
and the definition (UX.72) give both exact identities
(UX.68).  Notice that $R$ need not itself be a DVR;
its actual complete local residue-field property is
what the universal space requires.

Finally, a locally constant map from a compact Galois
group to the discrete ambient residue field has finite
image.  Hence every residual representation and its
finitely many flags use only finitely many residue
coefficients.  The same assertion holds for an actual
polynomial link: its matrix-entry maps into
$k_V[X]$ are locally constant, so only finitely many
polynomials, and therefore finitely many polynomial
coefficients, occur.  Apply the ambient residue
approximation property to these finite sets for
every point and every link of the finite chain.
Their finite compositum is the required threshold.
For every larger $\Lambda$, all objects descend
entry by entry to $k_\Lambda$ or $k_\Lambda[X]$.
Injectivity of $j_\Lambda$ descends multiplication,
symplectic, flag, inertia and support identities.
It also preserves nonzero-entry witnesses and the
nontrivial-character witnesses in goodness.  Since
the original maps have finite image, local constancy
descends as well.  Evaluation at $0$ and $1$ commutes
with the coefficient injection, so the actual links
have exactly their original endpoints.
\end{proof}

\begin{theorem}[UX12: the full case-II support witness]
Assume the literal \texttt{SeedIIAlt} and
\texttt{PMExAlt}.  For every original case-II datum
$D$, every fixed $\iota$, every ambient valued field
and every actual small \texttt{KPair} $(r,u)$ with
$r_K$ \texttt{Star}, the literal
\texttt{ProModAlt}$\,(D,\iota,V,r)$ holds.
The only further inputs are the cohomology and
Ferrero--Washington theorems already supplied by
BX and TW, and the residual vanishing supplied
by BR and AX.  In particular the original
representation, its multiplier, its ordered inertia
characters and the fixed $\iota$ are retained.
\end{theorem}

\begin{proof}
Choose the permitted label from the first-column
residual ordinary flag, as in the first paragraph
of the proof of UX2.  Apply UX1 to this label to
obtain the initial tame tower $T_0$ and its
separating element $g_0$.  Apply the remaining
construction of UX2 at this $g_0$ to put the
global representation in the normalized lattice
shape.  Let $r_1$ denote this integral representation
over $V$.  The change is conjugation over the
original $K$, and its multiplier is still
$a_O(D.\vartheta)$.  UX3--UX4 provide a later tame tower
$T_1$ such that every subsequent deep tower admits
an actual residual chain from the restriction of
this lattice to an actual integral seed.
All this is performed with the single initially
fixed $\iota$; the seed input has a universal
quantifier over $\iota$ and is applied at that
value only.

Apply UX5 at $T_1$ to choose its constant $C$.
Only afterwards choose a cyclotomic layer $T$
which is deep and satisfies
$d\ge30(C+1)^2$.  The bounds are uniform for
all residue coefficient fields and all later
layers, so they still hold at this $T$.
UX4 and the literal seed hypothesis now give
some $O_1\supset O$ and an actual integral
pair $Q$ over $O_1$ with
$T.\mathtt{SeedPairAlt}\,(\iota,Q)$.
UX11 gives the small restricted pair $P$ of
$r_1$, satisfying (UX.67).  The chain has the
form
\[
 \overline P\ \longleftrightarrow\ \cdots\
 \longleftrightarrow\ \overline Q
 \quad\hbox{in the one ambient residue field.}
 \tag{UX.73}
\]
It is the chain for this actual restriction and
these chosen places, not a chain to an isomorphic
representation on a newly chosen tower.

Take four coefficient thresholds: $O_1$ for the
integral seed, $O_2$ for the descent of (UX.73),
$O_3$ for the small-pair descent of UX11, and
$O_4$ from \texttt{PMExAlt} at this $D,T,\iota$.
Choose a single coefficient DVR $\Lambda$
containing them all.  Such a finite compositum
is again an allowed coefficient DVR inside the
fixed $\overline{\mathbb Q}_\ell$.
The order of these choices matters:
\texttt{PMExAlt} supplies $O_4$ before the
universal space is chosen, and promises its
conclusion for \emph{every} space over every
later coefficient DVR.  We therefore use UX6
to construct the actual space $X=(A,B,P_X)$
over $\Lambda$, and then obtain its ideal
$J\subset B$ with all five clauses of
$T.\mathtt{PotentialModularityTower}\,(\iota,X,J)$.
They include the stated finiteness, every
finite-field power-series point, every seed,
and both characteristic-zero conclusions.
No replacement of those clauses by a pointwise
or generic version has occurred.

Write $s,t$ for the descended endpoints of
(UX.73) in $k_\Lambda$.  UX11 gives an
injective local map $f:R\to V$ and a
continuous pair $P_R$ lifting $s$, with
$P_R\otimes_fV=P$.  The seed extends to
the integral pair $Q_\Lambda$ over
$\Lambda$, lifting $t$.
All equalities follow from the actual
coefficient inclusions and the injective
residue map $j_\Lambda$.  Universality of
this one space gives
\[
 f_0:A\longrightarrow R,
 \qquad f_1:A\longrightarrow\Lambda,
 \qquad P_X\otimes_{f_0}R=P_R,
 \quad P_X\otimes_{f_1}\Lambda=Q_\Lambda.
 \tag{UX.74}
\]
These are maps of $\Lambda$-algebras; their
reductions are precisely $s$ and $t$.

Compose $f_1$ with the given inclusion
$\Lambda\hookrightarrow\overline{\mathbb Q}_\ell$.
The resulting point is the original seed at
the fixed $\iota$, so the \texttt{seed}
clause gives \texttt{LocZero}.  Choose a
minimal prime $q$ of $A$ inside the kernel
of this characteristic-zero point.
Its contraction to $B$ is a prime annihilated
by the point; hence \texttt{LocZero} says
$J\subset q\cap B$.
The kernel is contained in the kernel of
the residue point $t$, since the map
$\Lambda\to\overline{\mathbb Q}_\ell$
is injective.  Thus
\[
 \mathtt{SomeIn}(J,\ker t).
 \tag{UX.75}
\]
This step does not assume that the seed
already annihilates $J$; support follows
from the stronger literal seed clause.

UX7--UX10 apply at the chosen uniform
bound and prove \texttt{AllIn} at every
good endpoint along (UX.73), in reverse
order starting from (UX.75).  In particular
every minimal prime $q\subset\ker s$
contains $JA$.
Set
\[
 y=f\circ f_0:A\longrightarrow V.
 \tag{UX.76}
\]
Its kernel is prime because $V$ is a
domain.  It lies in $\ker s$: injectivity
of $f$ says that $y(z)=0$ implies
$f_0(z)=0$, and then the reduction of
$f_0(z)$ is zero.  Choose a minimal prime
inside $\ker y$.  The just proved
\texttt{AllIn} statement gives
$JA\subset\ker y$.
Thus $y$ annihilates $J$ on the actual
valued point.

The remaining defining clauses can now
be checked as equalities rather than
existential compatibility assertions.
From (UX.68) and (UX.74),
\[
 \begin{gathered}
 y\circ a_{\Lambda,A}=a_\Lambda,
 \qquad y(\mathfrak m_B)\subset\mathfrak m_V,
 \qquad P_X\otimes_yV=P,\\
 r_1|_{G_F}=(P_X\otimes_yV).r,\qquad
 y\bigl(P_X.\psi_{v,i}(s)\bigr)
    =a_O\bigl(D.\mathrm{ps}_i(s|_{\mathbb Q})\bigr)
       \quad(s\in I_{F_v}).
 \end{gathered}
 \tag{UX.77}
\]
For the maximal-ideal assertion, the
restriction $B\to R$ is a map of
complete local $\Lambda$-algebras with
the same residue field; (UX.68) carries
its zero residues to zero residues in
$V$.  The last assertion is exactly
(UX.67), with the original indices.
The tower is fit and deep, its full
five-clause ideal was supplied above,
$r_1$ has the original multiplier, and
$r_K$ and $(r_1)_K$ are conjugate over
the same ambient $K$.  Together these
are every clause of the literal
\texttt{ProModAlt} definition.

All universal choices in the conclusion
were arbitrary.  The existential tower
and coefficient DVR may depend on the
original point, as the source permits;
the residue bounds are uniform before
their use, and the \texttt{PMExAlt}
coefficient threshold is uniform in
the subsequently chosen space.
No change of $\iota$ is made anywhere
in the construction.
\end{proof}

\subsection{The exact remaining boundary}

UX12 closes the original case-II implication
\texttt{proMod\_of\_pmEx\_of\_ctx2\_of\_kPair\_of\_star}.
It supplies the actual witness consumed by IX10,
including all coefficient, field, inertia and
support equalities.  Its two automorphic inputs
are the literal \texttt{SeedIIAlt} and
\texttt{PMExAlt}; their independent producers
belong to SX and QX.  Applying any such producer
must preserve its own hypotheses.  In particular
this chapter does not infer the five-clause
support ideal from generic automorphic finiteness.

The case-III analogue, based on
\texttt{SeedIVAlt}, has additional tower and
residual-pattern quantifiers.  UX12 makes no
claim to that analogue or to universal
\texttt{SymplecticFourImageDense}.  Consequently
the older continuation targets combining both
cases must retain their unsupplied case-III
scope even after this case-II chain is installed.

